\chapter{部分习题  答案与提示}

\section{第一章 实数集与函数}

\subsection{习题 1.1}

4. 当 \(x =  \pm  1\) 时等号成立.

9. (1) 当 \(a < b\) 时, \(x < \frac{a + b}{2}\) ; 当 \(a > b\) 时, \(x > \frac{a + b}{2}\) ;

(2)当 \(a > b\) 时, \(x > \frac{a + b}{2}\) ；

(3) 当 \(a \geq  b > 0\) 时, \(\sqrt{a - b} < \left| x\right|  < \sqrt{a + b}\) ; 当 \(\left| a\right|  < b\) 时, \(\left| x\right|  < \sqrt{a + b}\) .

\subsection{习题 1.2}

1. (1) \(x \in  \left( {-\infty ,\frac{1}{2}}\right\rbrack\) ;

(2) \(x \in  \left\lbrack  {-3 - 2\sqrt{2}, - 3 + 2\sqrt{2}}\right\rbrack   \cup  \left\lbrack  {3 - 2\sqrt{2},3 + 2\sqrt{2}}\right\rbrack\) ;

(3) \(x \in  \left( {a,b}\right)  \cup  \left( {c, + \infty }\right)\) ;

(4) \(x \in  \left\lbrack  {\frac{\pi }{4} + {2k\pi },\frac{3}{4}\pi  + {2k\pi }}\right\rbrack  ,k = 0, \pm  1, \pm  2,\cdots\) .

4. (1) sup \(S = \sqrt{2}\) , inf \(S =  - \sqrt{2}\) ; (2) sup \(S =  + \infty\) , inf \(S = 1\) ;

(3) \(\sup S = 1,\inf S = 0;\left( 4\right) \sup S = 1,\inf S = \frac{1}{2}\) .

\subsection{习题 1.3}

3. \({f}_{1}\left( x\right)  = \left\{  {\begin{matrix} {4x}, & 0 \leq  x \leq  \frac{1}{2}, \\  4 - {4x}, & \frac{1}{2} < x \leq  1; \end{matrix}\;{f}_{2}\left( x\right)  = \left\{  \begin{array}{ll} {16x}, & 0 \leq  x \leq  \frac{1}{4}, \\  8 - {16x}, & \frac{1}{4} < x \leq  \frac{1}{2}, \\  0, & \frac{1}{2} < x \leq  1. \end{array}\right. }\right.\)

4. (1) \(\left( {-\infty , + \infty }\right) ;\left( 2\right) \left( {1, + \infty }\right) ;\left( 3\right) \left\lbrack  {1,{100}}\right\rbrack  ;\left( 4\right) (0,{10}\rbrack\) .

5. (1) \(- 1,2,2;\left( 2\right) {2}^{\Delta x} - 2, - {\Delta x}\) .

6. \(\frac{1}{3 + x},\frac{1}{1 + {2x}},\frac{1}{1 + {x}^{2}},\frac{1 + x}{2 + x},\frac{1}{2 + x}\) .

7. (1) \(y = {u}^{20},u = 1 + x\) ; (2) \(y = {u}^{2},u = \arcsin v,v = {x}^{2}\) ;

(3) \(y = \lg u,u = 1 + v,v = \sqrt{w},w = 1 + {x}^{2}\) ; (4) \(y = {2}^{u},u = {v}^{2},v = \sin x\) .

10. (1) 成立; (2) 不成立.

\subsection{习题 1.4}

4. (1) 偶; (2) 奇; (3) 偶; (4) 奇.

5. (1) \(\pi ;\left( 2\right) \frac{\pi }{3};\left( 3\right) {12\pi }\) .

\subsection{第一章总练习题}

2. 是初等函数. (提示: 利用第 1 题的结果.)

3. \(\frac{1 + x}{1 - x}, - \frac{x}{2 + x},\frac{2}{1 + x},\frac{x - 1}{x + 1},\frac{1 + x}{1 - x},\frac{1 - {x}^{2}}{1 + {x}^{2}},x\) .

4. \(\frac{1}{x} + \frac{\sqrt{{x}^{2} + 1}}{\left| x\right| }\) .

5. (1) \(y = \left\lbrack  \frac{x + 2}{5}\right\rbrack  ,x = {30},{31},\cdots ,{50};\left( 2\right) y = \left\lbrack  {x + {0.5}}\right\rbrack  ,x > 0\) .

14. (1) (i) \(f\left( x\right)  = \left\{  {\begin{array}{ll} \sin x + 1, & x > 0, \\  0, & x = 0, \\  \sin x - 1, & x < 0, \end{array}\;\text{ (ii) }f\left( x\right)  = \left\{  \begin{array}{ll} \sin x + 1, & x \geq  0, \\  1 - \sin x, & x < 0; \end{array}\right. }\right.\)

(2) (i) \(f\left( x\right)  = \left\{  {\begin{array}{ll} {x}^{3}, & x <  - 1, \\   - 1 + \sqrt{1 - {x}^{2}}, &  - 1 \leq  x \leq  0, \\  1 - \sqrt{1 - {x}^{2}}, & 0 < x \leq  1, \\  {x}^{3}, & x > 1, \end{array}\text{ (ii) }f\left( x\right)  = \left\{  \begin{array}{ll}  - {x}^{3}, & x <  - 1, \\  1 - \sqrt{1 - {x}^{2}}, & \left| x\right|  \leq  1, \\  {x}^{3}, & x > 1. \end{array}\right. }\right.\)

\section{第二章 数 列 极 限}

\subsection{习题 2.1}

3. (1) 0,无穷小数列; (2) 1; (3) 0,无穷小数列;

(4) 0,无穷小数列；(5) 0,无穷小数列；(6) 1 ； (7) 1.

7. (1) 无界数列；(2) 有界数列；(3) 无穷大量；(4) 无界数列.

\subsection{习题 2.2}

1. (1) \(\frac{1}{4};\left( 2\right) 0;\left( 3\right) \frac{1}{3};\left( 4\right) \frac{1}{2};\left( 5\right) {10};\left( 6\right) 2\) .

4. (1) 1; (2) 2; (3) 3; (4) 1; (5) 0; (6) 1.

8. (1) 0 (提示: 先证明 \(\frac{1}{2} \cdot  \frac{3}{4} \cdot  \cdots  \cdot  \frac{{2n} - 1}{2n} < \frac{1}{\sqrt{{2n} + 1}}\) );

(2)1(提示: \(n! < \mathop{\sum }\limits_{{p = 1}}^{n}p! < \left( {n - 2}\right) \left( {n - 2}\right) ! + \left( {n - 1}\right) ! + n! < 2\left( {n - 1}\right) ! + n!\) );

(3)0(提示:先证明 \(0 < {\left( n + 1\right) }^{\alpha } - {n}^{\alpha } \leq  {n}^{\alpha  - 1}\) )；

(4) \(\frac{1}{1 - \alpha }\) (提示:记 \({p}_{n} = \left( {1 + \alpha }\right) {\left( 1 + \alpha \right) }^{2}\cdots {\left( 1 + \alpha \right) }^{{2}^{n}}\) ,则 \(\left( {1 - \alpha }\right) {p}_{n} = 1 - {\alpha }^{{2}^{n + 1}}\) ).

\subsection{习题 2.3}

1. (1) \(\frac{1}{\mathrm{e}};\left( 2\right) \mathrm{e};\left( 3\right) \mathrm{e};\left( 4\right) \sqrt{\mathrm{e}};\left( 5\right)\) 1.

3. (1) \(2;\left( 2\right) \frac{1}{2}\left( {1 + \sqrt{1 + {4c}}}\right) ;\left( 3\right) 0\) .

\subsection{第二章总练习题}

1. (1) 3; (2) 0; (3) 0.

12. 否, 反例如:

\[
\left\{  {a}_{n}\right\}   = 1,\frac{1}{2},3,\frac{1}{4},5,\frac{1}{6},\cdots ,
\]

\[
\left\{  {b}_{n}\right\}   = 1,2,\frac{1}{3},4,\frac{1}{5},6\cdots ,
\]

\[
\left\{  {{a}_{n}{b}_{n}}\right\}   = 1,1,1,1,1,1,\cdots .
\]

\section{第三章 函 数 极 限}

\subsection{习题 3.1}

6. (1) \(f\left( {0 - 0}\right)  =  - 1,f\left( {0 + 0}\right)  = 1\) ; (2) \(f\left( {0 - 0}\right)  =  - 1,f\left( {0 + 0}\right)  = 0\) ;

(3) \(f\left( {0 - 0}\right)  = f\left( {0 + 0}\right)  = 1\) .

\subsection{习题 3.2}

1. (1) \(2 - \frac{{\pi }^{2}}{2}\) ; (2)1; (3) \(\frac{2}{3}\) ; (4)-3; (5) \(\frac{n}{m}\) ; (6) \(\frac{4}{3}\) ; (7) \(\frac{1}{2a}\) ; (8) \(\frac{{3}^{70}{8}^{20}}{{5}^{90}}\) .

2. (1) 1; (2) 0 .

4. \(m < n\) 时,0; \(m = n\) 时, \(\frac{{a}_{0}}{{b}_{0}}\) .

8. (1) \(- 1;\left( 2\right) 1;\left( 3\right)  - 3;\left( 4\right) \frac{1}{n};\left( 5\right) 1\) .

\subsection{习题 3.4}

1. (1) 2; (2) 0; (3) -1; (4) 1; (5) 2; (6) 1; (7) 1; (8) sin2a; (9) 8; (10) \(\sqrt{2}\) .

2. (1) \({\mathrm{e}}^{2}\) ; (2) \({\mathrm{e}}^{\alpha }\) ; (3) e; (4) \({\mathrm{e}}^{2}\) ; (5) \({\mathrm{e}}^{2}\) ; (6) \({\mathrm{e}}^{\alpha \beta }\) .

4. (1) \(0;\left( 2\right) \mathrm{e}\) .

\subsection{习题 3.5}

2. (1) 0 ; (2) 1 .

4. (1) \(y = 0,x = 0;\left( 2\right) y = \frac{\pi }{2},y =  - \frac{\pi }{2};\left( 3\right) y = {3x} + 6,x = 0,x = 2\) .

5. (1) 3; (2) 2; (3) 1; (4) \(\frac{2}{5}\) .

6. (1) \(\frac{5}{2}\) ; (2) 2; (3) \(\frac{1}{2}n\left( {n + 1}\right)\) .

\subsection{第三章总练习题}

1. (1) \(1;\left( 2\right) \frac{1}{2};\left( 3\right) a + b\) ;

(4) \(1;\left( 5\right)  - 1;\left( 6\right) \frac{3}{2};\left( 7\right) \frac{1}{2}\left( {m - n}\right)\) .

2. (1) \(a = 1,b =  - 1;\left( 2\right) a =  - 1,b = \frac{1}{2};\left( 3\right) a = 1,b =  - \frac{1}{2}\) .

8. (1) \(+ \infty ;\left( 2\right) 0\) .

10. (1) \(+ \infty ;\left( 2\right)  + \infty\) .

\section{第四章 函数的连续性}

\subsection{习题 4.1}

2. (1) \(x = 0\) ,第二类间断点; (2) \(x = 0\) ,跳跃间断点;

(3) \(x = {n\pi }\left( {n = 0, \pm  1, \pm  2,\cdots }\right)\) ,可去间断点; (4) \(x = 0\) ,可去间断点;

(5) \(x = \frac{\pi }{2} + {k\pi }\left( {k = 0, \pm  1, \pm  2,\cdots }\right)\) ,跳跃间断点; (6) 除 \(x = 0\) 外每一点都是第二类间断点;

(7) \(x =  - 7\) 为第二类间断点, \(x = 1\) 为跳跃间断点.

\subsection{习题 4.2}

1. (1) \(f \circ  g\) 处处连续, \(g \circ  f,x = 0\) 为可去间断点;

(2) \(f \circ  g,x =  - 1,0,1\) 为跳跃间断点, \(g \circ  f\) 处处连续.

8. (1) \(\frac{3}{4}\pi ;\left( 2\right) \frac{\sqrt{3}}{2}\) .

\subsection{习题 4.3}

1. (1) 6; (2) \(\frac{1}{2}\) ; (3) 1; (4) 1; (5) e.

\section{第五章 导数和微分}

\subsection{习题 5.1}

1. \({\Delta t} = 1,\bar{v} = {55};{\Delta t} = {0.1},\bar{v} = {50.5};{\Delta t} = {0.01},\bar{v} = {50.05};v = {50}\) .

2. 设 \(t\) 时刻所对应的旋转角为 \(\Phi \left( t\right)\) ,则角速度为 \(\omega \left( t\right)  = \frac{\mathrm{d}\Phi \left( t\right) }{\mathrm{d}t}\) .

3. 4.

4. \(a = 6,b =  - 9\) .

5. (1) \(\left( {1,0}\right) ;\left( 2\right) \left( {\frac{1}{2}, - \ln 2}\right)\) .

6. (1) 切线方程: \(y = x - 1\) ,法线方程: \(y =  - x + 3\) ;

(2)切线方程: \(y = 1\) ,法线方程: \(x = 0\) .

7. (1) \({f}^{\prime }\left( x\right)  = \left\{  {\begin{matrix} 3{x}^{2}, & x \geq  0, \\   & \left( 2\right) {f}^{\prime }\left( x\right)  = \left\{  \begin{matrix} 1, & x > 0, \\  \text{ 不存在, } & x = 0, \\  0, & x < 0; \end{matrix}\right.  \end{matrix}\text{ , }}\right.\)

8. (1) \(m \geq  1;\) (2) \(m \geq  2\) .

9. (1) \({k\pi } - \frac{\pi }{4}\left( {k = 0, \pm  1, \pm  2,\cdots }\right) ;\;\left( 2\right) x = 1\) .

11. 0 .

15. \(\frac{\pi }{2} - \arctan \frac{2}{5}\) .

\subsection{习题 5.2}

1. (1) \({f}^{\prime }\left( 0\right)  = 0,{f}^{\prime }\left( 1\right)  = {18};\left( 2\right) {f}^{\prime }\left( 0\right)  = 1,{f}^{\prime }\left( \pi \right)  =  - 1\) ;

(3) \({f}^{\prime }\left( 1\right)  = \frac{1}{4\sqrt{2}},{f}^{\prime }\left( 4\right)  = \frac{1}{8\sqrt{3}}\) .

2. (1) \({y}^{\prime } = {6x}\) ; (2) \({y}^{\prime } = \frac{-{x}^{2} - {4x} - 1}{{\left( {x}^{2} + x + 1\right) }^{2}}\) ; (3) \({y}^{\prime } = n\left( {{x}^{n - 1} + 1}\right)\) ;

(4) \({y}^{\prime } = \frac{1}{m} - \frac{m}{{x}^{2}} + \frac{1}{\sqrt{x}} - \frac{1}{x\sqrt{x}};\;\left( 5\right) {y}^{\prime } = 3{x}^{2}{\log }_{3}x + \frac{{x}^{2}}{\ln 3}\) ; (6) \({y}^{\prime } = {\mathrm{e}}^{x}\left( {\cos x - \sin x}\right)\) ; (7) \({y}^{\prime } =  - {18}{x}^{5} + 5{x}^{4} - {12}{x}^{3} + {12}{x}^{2} - {2x} + 3\) ; (8) \({y}^{\prime } = \frac{x{\sec }^{2}x - \tan x}{{x}^{2}}\) ; (9) \(\frac{1 - \cos x - x\sin x}{{\left( 1 - \cos x\right) }^{2}}\) ;

(10) \({y}^{\prime } = \frac{2}{x{\left( 1 - \ln x\right) }^{2}}\) ; (11) \(\frac{1}{2\sqrt{x}}\arctan x + \frac{\sqrt{x} + 1}{1 + {x}^{2}}\) ;

(12) \({y}^{\prime } = \frac{{2x}\left( {\sin x + \cos x}\right)  - \left( {{x}^{2} + 1}\right) \left( {\cos x - \sin x}\right) }{{\left( \sin x + \cos x\right) }^{2}}\) .

3. (1) \({y}^{\prime } = \frac{1 - 2{x}^{2}}{\sqrt{1 - {x}^{2}}}\) ; (2) \({y}^{\prime } = {6x}{\left( {x}^{2} - 1\right) }^{2}\) ;

(3) \({y}^{\prime } = 3 \cdot  \frac{{\left( 1 + {x}^{2}\right) }^{2}\left( {1 + {2x} - {x}^{2}}\right) }{{\left( 1 - x\right) }^{4}}\) ; (4) \({y}^{\prime } = \frac{1}{x\ln x}\) ;

(5) \({y}^{\prime } = \cot x\) ; (6) \({y}^{\prime } = \frac{{2x} + 1}{{x}^{2} + x + 1} \cdot  \frac{1}{\ln {10}}\) ;

(7) \({y}^{\prime } = \frac{1}{\sqrt{1 + {x}^{2}}}\) ; (8) \({y}^{\prime } = \frac{1}{x\sqrt{1 - {x}^{2}}}\) ;

(9) \({y}^{\prime } = 3\cos {2x}\left( {\sin x + \cos x}\right)\) ; (10) \({y}^{\prime } =  - 6\cos {4x}\sin {8x}\) ;

(11) \({y}^{\prime } = \frac{x}{\sqrt{1 + {x}^{2}}}\cos \sqrt{1 + {x}^{2}}\) ; (12) \({y}^{\prime } = {6x}{\sin }^{2}{x}^{2}\cos {x}^{2}\) ;

(13) \({y}^{\prime } = \frac{-1}{\left| x\right| \sqrt{{x}^{2} - 1}}\) ; (14) \({y}^{\prime } = \frac{6{x}^{2}}{1 + {x}^{6}}\arctan {x}^{3}\) ;

(15) \({y}^{\prime } = \frac{-1}{1 + {x}^{2}}\) ; (16) \({y}^{\prime } = \frac{\sin {2x}}{\sqrt{1 - {\sin }^{4}x}}\) ;

(17) \({y}^{\prime } = {\mathrm{e}}^{x + 1}\) ; (18) \({y}^{\prime } = \ln 2 \cdot  {2}^{\sin x}\cos x\) ;

(19) \({y}^{\prime } = {x}^{\sin x}\left( {\cos x\ln x + \frac{\sin x}{x}}\right)\) ; (20) \({y}^{\prime } = {x}^{{x}^{x}}{x}^{x}\left\lbrack  {{\ln }^{2}x + \ln x + \frac{1}{x}}\right\rbrack\) ;

(21) \({y}^{\prime } = {\mathrm{e}}^{-x}\left\lbrack  {2\cos {2x} - \sin {2x}}\right\rbrack  ;\left( {22}\right) {y}^{\prime } = \frac{4\sqrt{x}\sqrt{x + \sqrt{x}} + 2\sqrt{x} + 1}{8\sqrt{x}\sqrt{x + \sqrt{x}} \cdot  \sqrt{x + \sqrt{x + \sqrt{x}}}}\) ;

(23) \({y}^{\prime } = \cos \left( {\sin \left( {\sin x}\right) }\right)  \cdot  \cos \left( {\sin x}\right)  \cdot  \cos x\) ;

(24) \({y}^{\prime } = \cos \left( \frac{x}{\sin \left( \frac{x}{\sin x}\right) }\right) \frac{\sin \left( \frac{x}{\sin x}\right)  - x\cos \left( \frac{x}{\sin x}\right) \frac{\sin x - x\cos x}{{\sin }^{2}x}}{{\sin }^{2}\left( \frac{x}{\sin x}\right) }\) ;

(25) \({y}^{\prime } = \left( {\mathop{\prod }\limits_{{j = 1}}^{n}{\left( x - {a}_{j}\right) }^{{\alpha }_{j}}}\right) \left( {\mathop{\sum }\limits_{{k = 1}}^{n}\frac{{\alpha }_{k}}{x - {a}_{k}}}\right)\) ;

(26) \({y}^{\prime } = \frac{\cos x}{\left| {a + b\sin x}\right| \left| {\cos x}\right| }\) .

4. (1) \({f}^{\prime }\left( x\right)  = 3{x}^{2},{f}^{\prime }\left( {x + 1}\right)  = 3{\left( x + 1\right) }^{2},{f}^{\prime }\left( {x - 1}\right)  = 3{\left( x - 1\right) }^{2}\) ;

(2) \({f}^{\prime }\left( x\right)  = 3{\left( x - 1\right) }^{2},{f}^{\prime }\left( {x + 1}\right)  = 3{x}^{2},{f}^{\prime }\left( {x - 1}\right)  = 3{\left( x - 2\right) }^{2}\) ;

(3) \({f}^{\prime }\left( x\right)  = 3{\left( x + 1\right) }^{2},{f}^{\prime }\left( {x + 1}\right)  = 3{\left( x + 2\right) }^{2},{f}^{\prime }\left( {x - 1}\right)  = 3{x}^{2}\) .

5. (1) \({f}^{\prime }\left( x\right)  = {g}^{\prime }\left( {x + g\left( a\right) }\right) ;\;\left( 2\right) {f}^{\prime }\left( x\right)  = {g}^{\prime }\left( {x + g\left( x\right) }\right) \left( {1 + {g}^{\prime }\left( x\right) }\right)\) ;

(3) \({f}^{\prime }\left( x\right)  = {g}^{\prime }\left( {{xg}\left( a\right) }\right)  \cdot  g\left( a\right)\) ;

(4) \({f}^{\prime }\left( x\right)  = {g}^{\prime }\left( {{xg}\left( x\right) }\right) \left( {g\left( x\right)  + x{g}^{\prime }\left( x\right) }\right)\) .

8. (1) \({y}^{\prime } = 3{\sinh }^{2}x\cosh x\) ; (2) \({y}^{\prime } = \sinh \left( {\sinh x}\right) \cosh x\) ;

(3) \({y}^{\prime } = \tanh x\) ; (4) \({y}^{\prime } = \frac{1}{{\sinh }^{2}x + {\cosh }^{2}x}\) .

9. (1) \({y}^{\prime } = \frac{1}{\sqrt{1 + {x}^{2}}}\) ; (2) \({y}^{\prime } = \frac{1}{\sqrt{{x}^{2} - 1}}\) ;

(3) \({y}^{\prime } = \frac{1}{1 - {x}^{2}}\left( {\left| x\right|  < 1}\right)\) ; (4) \({y}^{\prime } = \frac{1}{1 - {x}^{2}}\left( {\left| x\right|  > 1}\right)\) ;

(5) \({y}^{\prime } = 0\) ; (6) \({y}^{\prime } = \left| {\sec x}\right|\) .

\subsection{习题 5.3}

1. (1) \({\left. \frac{\mathrm{d}y}{\mathrm{\;d}x}\right| }_{t = \frac{\pi }{3}} =  - 3;\left( 2\right) \frac{\mathrm{d}y}{\mathrm{\;d}x} =  - 2\) .

2. (1) \({\left. \frac{\mathrm{d}y}{\mathrm{\;d}x}\right| }_{t = \frac{\pi }{2}} = 1,{\left. \;\frac{\mathrm{d}y}{\mathrm{\;d}x}\right| }_{t = \pi } = 0\) .

3. (1) 切线方程 \(y = \frac{1}{2}x\) ,法线方程 \(y =  - {2x}\) ;

(2) \({2y} - \left( {2 - \sqrt{2}}\right) x = \frac{3}{2}\sqrt{2} - 2,\;{2x} + \left( {2 - \sqrt{2}}\right) y = \frac{3}{2}\sqrt{2} - 1\) .

6. \(\frac{1}{2}\left( {\theta  - \pi }\right)\) .

\subsection{习题 5.4}

\(1\;\left( 1\right) {y}^{\prime \prime }\left( 1\right)  = {26},{y}^{\prime \prime \prime }\left( 1\right)  = {18},{y}^{\left( 4\right) }\left( 1\right)  = 0;\)

(2) \({f}^{\prime \prime }\left( 0\right)  = 0,{f}^{\prime \prime }\left( 1\right)  =  - \frac{3}{4\sqrt{2}},{f}^{\prime \prime }\left( {-1}\right)  = \frac{3}{4\sqrt{2}}\) .

3. (1) \({f}^{\prime \prime }\left( x\right)  = \frac{1}{x};\;\left( 2\right) {f}^{\prime \prime \prime }\left( x\right)  = {4x}{\mathrm{e}}^{-{x}^{2}}\left( {3 - 2{x}^{2}}\right)\) ;

(3) \({f}^{\left( 5\right) }\left( x\right)  = \frac{24}{{\left( 1 + x\right) }^{5}}\) ;

4. (1) \({y}^{\prime \prime } = \frac{1}{{x}^{2}}\left\lbrack  {{f}^{\prime \prime }\left( {\ln x}\right)  - {f}^{\prime }\left( {\ln x}\right) }\right\rbrack\) ;

(2) \({y}^{\prime \prime } = n\left( {n - 1}\right) {x}^{n - 2}{f}^{\prime }\left( {x}^{n}\right)  + {\left( n{x}^{n - 1}\right) }^{2}{f}^{\prime \prime }\left( {x}^{n}\right)\) ;

(3) \({y}^{\prime \prime } = {f}^{\prime \prime }\left( {f\left( x\right) }\right) {\left( {f}^{\prime }\left( x\right) \right) }^{2} + {f}^{\prime }\left( {f\left( x\right) }\right) {f}^{\prime \prime }\left( x\right)\) .

5. (1) \({y}^{\left( n\right) } = \frac{{\left( -1\right) }^{n - 1}\left( {n - 1}\right) !}{{x}^{n}};\;\left( 2\right) {y}^{\left( n\right) } = {a}^{x}{\ln }^{n}a\) ;

(3) \({y}^{\left( n\right) } = n!\left( {\frac{{\left( -1\right) }^{n}}{{x}^{n + 1}} + \frac{1}{{\left( 1 - x\right) }^{n + 1}}}\right)\) ;

(4) \({y}^{\left( n\right) } = \frac{{\left( -1\right) }^{n}n!\left( {\ln x - \mathop{\sum }\limits_{{k = 1}}^{n}\frac{1}{k}}\right) }{{x}^{n + 1}}\) ;

(5) \({y}^{\left( n\right) } = \frac{n!}{{\left( 1 - x\right) }^{n + 1}}\) ;

(6) \({y}^{\left( n\right) } = {\left( {a}^{2} + {b}^{2}\right) }^{\frac{n}{2}}{\mathrm{e}}^{ax}\sin \left( {{bx} + n{\varphi }_{0}}\right) ,{\varphi }_{0} = \arctan \frac{b}{a}\) .

6. (1) \({y}^{\prime \prime } = \frac{1}{{3a}\sin t{\cos }^{4}t};\left( 2\right) {y}^{\prime \prime } = \frac{2}{{\mathrm{e}}^{t}{\left( \cos t - \sin t\right) }^{3}}\) .

7. \({f}^{\prime }\left( x\right)  = \left\{  {\begin{matrix} 3{x}^{2}, & x \geq  0, \\   - 3{x}^{2}, & x < 0, \end{matrix}\;{f}^{\prime \prime }\left( x\right)  = \left\{  \begin{matrix} {6x}, & x \geq  0, \\   - {6x}, & x < 0, \end{matrix}\right. }\right.\)

\({f}^{\prime \prime \prime }\left( x\right)  = \left\{  \begin{matrix} 6, & x > 0, \\  \text{ 不存在, } & x = 0, \\   - 6, & x < 0. \end{matrix}\right.\)

8. \({\left( {f}^{-1}\right) }^{\prime \prime \prime }\left( y\right)  = \frac{3{f}^{\prime \prime }{\left( x\right) }^{2} - {f}^{\prime }\left( x\right) {f}^{\prime \prime \prime }\left( x\right) }{{f}^{\prime }{\left( x\right) }^{5}}\) .

9. (2) \({\left. {y}^{\left( 2m\right) }\right| }_{x = 0} = 0,{\left. \;{y}^{\left( 2m + 1\right) }\right| }_{x = 0} = {\left( -1\right) }^{m}\left( {2m}\right)\) !.

10. (2) \({\left. {y}^{\left( 2m\right) }\right| }_{x = 0} = 0,{\left. \;{y}^{\left( 2m + 1\right) }\right| }_{x = 0} = {\left\lbrack  \left( 2m - 1\right) !!\right\rbrack  }^{2}\) .

\subsection{习题 5.5}

1. 当 \({\Delta x} = {0.1},{\Delta y} = {0.21},\mathrm{\;d}y = {0.2},{\Delta y} - \mathrm{d}y = {0.01}\) ; 当 \({\Delta x} = {0.01},{\Delta y} = {0.0201},\mathrm{\;d}y = {0.02},{\Delta y} - \mathrm{d}y =\) 0.000 1 .

2. (1) \(\mathrm{d}y = \left( {1 + {4x} - {x}^{2} + 4{x}^{3}}\right) \mathrm{d}x\) ; (2) \(\mathrm{d}y = \ln x\mathrm{\;d}x\) ;

(3) \(\mathrm{d}y = \left( {{2x}\cos {2x} - 2{x}^{2}\sin {2x}}\right) \mathrm{d}x\) ; (4) \(\mathrm{d}y = \frac{1 + {x}^{2}}{{\left( 1 - {x}^{2}\right) }^{2}}\mathrm{\;d}x\) ;

(5) \(\mathrm{d}y = {\mathrm{e}}^{ax}\left( {a\sin {bx} + b\cos {bx}}\right) \mathrm{d}x\) ; (6) \(\mathrm{d}y =  - \operatorname{sgn}x\frac{\mathrm{d}x}{\sqrt{1 - {x}^{2}}}\) .

4. (1) 1.007; (2) 0.9933; (3) 1.0058; (4) 5.1.

5. \({0.33}\%\) .

6. 弦长.

\subsection{第五章总练习题}

6. \({f}_{ + }^{\prime }\left( a\right)  = \varphi \left( a\right) ,{f}_{ - }^{\prime }\left( a\right)  =  - \varphi \left( a\right)\) ,当 \(\varphi \left( a\right)  = 0,{f}^{\prime }\left( a\right)\) 存在且等于零.

7. (1) \({y}^{\prime } = {\mathrm{e}}^{f\left( x\right) }\left( {{f}^{\prime }\left( {\mathrm{e}}^{x}\right) {\mathrm{e}}^{x} + f\left( {\mathrm{e}}^{x}\right) {f}^{\prime }\left( x\right) }\right)\) ;

(2) \({y}^{\prime } = {f}^{\prime }\left( {f\left( {f\left( x\right) }\right) }\right) {f}^{\prime }\left( {f\left( x\right) }\right) {f}^{\prime }\left( x\right)\) .

8. (1) \({y}^{\prime } = \frac{{\varphi }^{\prime }\left( x\right) \varphi \left( x\right)  + {\psi }^{\prime }\left( x\right) \psi \left( x\right) }{\sqrt{{\left( \varphi \left( x\right) \right) }^{2} + {\left( \psi \left( x\right) \right) }^{2}}}\left( {{\varphi }^{2}\left( x\right)  + {\psi }^{2}\left( x\right)  \neq  0}\right)\) ;

(2) \({y}^{\prime } = \frac{{\varphi }^{\prime }\left( x\right) \psi \left( x\right)  - \varphi \left( x\right) {\psi }^{\prime }\left( x\right) }{{\varphi }^{2}\left( x\right)  + {\psi }^{2}\left( x\right) }\) ;

(3) \({y}^{\prime } = \frac{{\psi }^{\prime }\left( x\right) \varphi \left( x\right) \ln \varphi \left( x\right)  - {\varphi }^{\prime }\left( x\right) \psi \left( x\right) \ln \psi \left( x\right) }{\varphi \left( x\right) \psi \left( x\right) {\ln }^{2}\varphi \left( x\right) }\) .

9. (1) \({F}^{\prime }\left( x\right)  = 3\left( {{x}^{2} + 5}\right) ;\;\left( 2\right) {F}^{\prime }\left( x\right)  = 6{x}^{2}\) .

\section{第六章 微分中值定理及其应用}

\subsection{习题 6.2}

5. (1) 1 ; (2) \(\frac{\sqrt{3}}{3}\) ; (3) 1 ; (4) 2 ; (5) 1 ; (6) \(\frac{1}{2}\) ;

(7) \(1;\left( 8\right) \frac{1}{\mathrm{e}};\left( 9\right) 1;\left( {10}\right) 0;\left( {11}\right)  - \frac{1}{3};\left( {12}\right) {\mathrm{e}}^{\frac{1}{3}}\) .

7. (1) \(- \frac{4}{{\pi }^{2}};\left( 2\right) 0;\left( 3\right) 1;\left( 4\right) {\mathrm{e}}^{-1}\) ;

(5) \(\frac{1}{2}\) ; (6)0; (7) \(- \frac{\mathrm{e}}{2}\) ; (8) \({\mathrm{e}}^{-1}\) .

\subsection{习题 6.3}

1. (1) \(f\left( x\right)  = 1 - \frac{1}{2}x + \frac{3!!}{2!{2}^{2}}{x}^{2} + \cdots  + {\left( -1\right) }^{n}\frac{\left( {{2n} - 1}\right) !!}{n!{2}^{n}}{x}^{n} + o\left( {x}^{n}\right)\) ;

(2) \(f\left( x\right)  = x - \frac{1}{3}{x}^{3} + \frac{1}{5}{x}^{5} + o\left( {x}^{5}\right)\) ;

(3) \(f\left( x\right)  = x + \frac{1}{3}{x}^{3} + \frac{2}{15}{x}^{5} + o\left( {x}^{5}\right)\) .

2. (1) \(\frac{1}{3};\;\left( 2\right) \frac{1}{2};\;\left( 3\right) \frac{1}{3}\) .

3. (1) \(f\left( x\right)  = {10} + {11}\left( {x - 1}\right)  + 7{\left( x - 1\right) }^{2} + {\left( x - 1\right) }^{3}\) ;

(2) \(f\left( x\right)  = 1 - x + {x}^{2} + \cdots  + {\left( -1\right) }^{n}{x}^{n} + \frac{{\left( -1\right) }^{n + 1}{x}^{n + 1}}{{\left( 1 + \theta x\right) }^{n + 2}},0 < \theta  < 1\) .

4. (1) \(\left| {{R}_{4}\left( x\right) }\right|  \leq  \frac{1}{{2}^{5} \cdot  5!}\) ; (2) \(\left| {{R}_{2}\left( x\right) }\right|  \leq  \frac{1}{16}\) .

5. (1) 取 \(n = {12},\mathrm{e} \approx  {2.718281828};\left( 2\right) {0.99325}\) .

\subsection{习题 6.4}

1. (1) 极大值 \(f\left( \frac{3}{2}\right)  = \frac{27}{16}\) ; (2)极小值 \(f\left( {-1}\right)  =  - 1\) ,极大值 \(f\left( 1\right)  = 1\) ;

(3)极小值 \(f\left( 1\right)  = 0\) ,极大值 \(f\left( {\mathrm{e}}^{2}\right)  = \frac{4}{{\mathrm{e}}^{2}}\) ；

(4)极大值 \(f\left( 1\right)  = \frac{\pi }{4} - \frac{1}{2}\ln 2\) .

4. (1) 最小值 \(f\left( {-1}\right)  =  - {10}\) ,最大值 \(f\left( 1\right)  = 2\) ;

(2)最大值 \(f\left( \frac{\pi }{4}\right)  = 1\) ,无最小值；

(3) 最小值 \(f\left( {\mathrm{e}}^{-2}\right)  =  - \frac{2}{\mathrm{e}}\) .

6. 边长为 \(\frac{l}{2}\) .

7. 半径与高之比为 \(1 : 1\) .

8. 取 \(x = \frac{{a}_{1} + {a}_{2} + \cdots  + {a}_{n}}{n}\) .

9. 取 \(a = 1\) .

10. (1) 极小值 \(f\left( 0\right)  = f\left( {\pm 1}\right)  = 0\) ,极大值 \(f\left( {\pm \frac{1}{\sqrt{3}}}\right)  = \frac{2}{3\sqrt{3}}\) ;

(2)极大值 \(f\left( 1\right)  = 2\) ,极小值 \(f\left( {-1}\right)  =  - 2\) ；

(3)极小值 \(f\left( 1\right)  = 0\) ,极大值 \(f\left( \frac{1}{5}\right)  = \frac{3456}{3125}\) .

11. \(a =  - \frac{2}{3},b =  - \frac{1}{6},{x}_{1}\) 极小值点, \({x}_{2}\) 极大值点.

12. \(\left( {p, \pm  \sqrt{2}p}\right)\) .

13. \(\frac{a\alpha }{\sqrt{{\beta }^{2} - {\alpha }^{2}}}\) .

\subsection{习题 6.5}

I. (1) 凹区间 \(\left( {-\infty ,\frac{1}{2}}\right)\) ,凸区间 \(\left( {\frac{1}{2}, + \infty }\right)\) ,拐点 \(\left( {\frac{1}{2},\frac{13}{2}}\right)\) ;

(2)凹区间 \(\left( {-\infty ,0}\right)\) ,凸区间 \(\left( {0, + \infty }\right)\) ；

(3)凹区间(-1,0),凸区间 \(\left( {-\infty , - 1}\right) ,\left( {0, + \infty }\right)\) ,拐点(-1,0)；

(4)凹区间 \(\left( {-\infty , - 1}\right) ,\left( {1, + \infty }\right)\) ,凸区间(-1,1),拐点 \(\left( {\pm 1,\ln 2}\right)\) ；

( 5 )凹区间 \(\left( {\frac{-1}{\sqrt{3}},\frac{1}{\sqrt{3}}}\right)\) ,凸区间 \(\left( {-\infty , - \frac{1}{\sqrt{3}}}\right) ,\left( {\frac{1}{\sqrt{3}}, + \infty }\right)\) ,拐点 \(\left( {\pm \frac{1}{\sqrt{3}},\frac{3}{4}}\right)\) .

2. \(a =  - \frac{3}{2},b = \frac{9}{2}\) .

习题 6.6

(1) III 部分习题答案与提示

\begin{center}
\adjustbox{max width=\textwidth}{
\begin{tabular}{|c|c|c|c|c|c|c|c|}
\hline
\(x\) & \(\left( {-\infty , - 5}\right)\) & -5 & (-5,-2) & -2 & (-2,1) & 1 & \(\left( {1, + \infty }\right)\) \\
\cline{1-8}
\({y}^{\prime }\) & + & 0 & - & - & - & 0 & + \\
\cline{1-8}
\({y}^{\prime \prime }\) & - & - & - & 0 & + & + & + \\
\cline{1-8}
\(y\) & 增凹 ↗ & 极大值 \(f\left( {-5}\right)  = {80}\) & 减凹 ↘ & 拐点 (-2,26) & 减凸 ↘ & 极小值 \(f\left( 1\right)  =  - {28}\) & 增凸 ↗ \\
\cline{1-8}
\hline
\end{tabular}
}
\end{center}

(2)

\begin{center}
\adjustbox{max width=\textwidth}{
\begin{tabular}{|c|c|c|c|c|c|c|}
\hline
\(x\) & \(\left( {-\infty , - 3}\right)\) & -3 & (-3,-1) & (-1,0) & 0 & \(\left( {0, + \infty }\right)\) \\
\cline{1-7}
\({y}^{\prime }\) & + & 0 & - & + & 0 & + \\
\cline{1-7}
\({y}^{\prime \prime }\) & - & - & - & - & 0 & + \\
\cline{1-7}
\(y\) & 增凹 ↗ & 极大值 \(f\left( {-3}\right)  =  - \frac{27}{8}\) & 减凹 ↘ & 增凹 ↗ & 拐点 (0,0) & 增凸 ↗ \\
\cline{1-7}
\hline
\end{tabular}
}
\end{center}

渐近线 \(x =  - 1,y = \frac{1}{2}x - 1\) ;

(3)

\begin{center}
\adjustbox{max width=\textwidth}{
\begin{tabular}{|c|c|c|c|c|c|c|c|}
\hline
\(x\) & \(\left( {-\infty , - 1}\right)\) & -1 & (-1,0) & 0 & (0,1) & 1 & \(\left( {1, + \infty }\right)\) \\
\cline{1-8}
\({y}^{\prime }\) & + & 0 & - & - & - & 0 & + \\
\cline{1-8}
\({y}^{\prime \prime }\) & - & - & - & 0 & + & + & + \\
\cline{1-8}
\(y\) & 增凹 ↗ & 极大值 \(f\left( {-1}\right)  =  - 1 + \frac{\pi }{2}\) & 减凹 ↘ & 拐点 (0,0) & 减凸 ↘ & 极小值 \(f\left( 1\right)  = 1 - \frac{\pi }{2}\) & 增凸 ↗ \\
\cline{1-8}
\hline
\end{tabular}
}
\end{center}

渐近线 \(y = x - \pi ,y = x + \pi\) ;

(4)

\begin{center}
\adjustbox{max width=\textwidth}{
\begin{tabular}{|c|c|c|c|c|c|}
\hline
\(x\) & \(\left( {-\infty ,1}\right)\) & 1 & (1,2) & 2 & \(\left( {2, + \infty }\right)\) \\
\cline{1-6}
\({y}^{\prime }\) & + & 0 & - & - & - \\
\cline{1-6}
\({y}^{\prime \prime }\) & - & - & - & 0 & + \\
\cline{1-6}
\(y\) & 增凹 ↗ & 极大值 \(f\left( 1\right)  = \frac{1}{\mathrm{e}}\) & 减凹 ↘ & 拐点 \(\left( {2,\frac{2}{\mathrm{e}}}\right)\) & 减凸 ↘ \\
\cline{1-6}
\hline
\end{tabular}
}
\end{center}

渐近线 \(y = 0\) ;

(5)奇函数

\begin{center}
\adjustbox{max width=\textwidth}{
\begin{tabular}{|c|c|c|c|c|c|c|}
\hline
\(x\) & 0 & \(\left( {0,\frac{1}{\sqrt{2}}}\right)\) & \(\frac{1}{\sqrt{2}}\) & \(\left( {\frac{1}{\sqrt{2}},1}\right)\) & 1 & \(\left( {1, + \infty }\right)\) \\
\cline{1-7}
\({y}^{\prime }\) & 0 & - & - & - & 0 & + \\
\cline{1-7}
\({y}^{\prime \prime }\) & 0 & - & 0 & + & + & + \\
\cline{1-7}
\(y\) & 拐点 (0,0) & 减凹 ↘ & 拐点 \(\left( {\frac{1}{\sqrt{2}}, - \frac{7}{4\sqrt{2}}}\right)\) & 减凸 ↘ & 极小值 \(f\left( 1\right)  =  - 2\) & 增凸 ↗ \\
\cline{1-7}
\hline
\end{tabular}
}
\end{center}

(6)偶函数

\begin{center}
\adjustbox{max width=\textwidth}{
\begin{tabular}{|c|c|c|c|c|}
\hline
\(x\) & 0 & \(\left( {0,\frac{1}{\sqrt{2}}}\right)\) & \(\frac{1}{\sqrt{2}}\) & \(\left( {\frac{1}{\sqrt{2}}, + \infty }\right)\) \\
\cline{1-5}
\({y}^{\prime }\) & 0 & - & - & - \\
\cline{1-5}
\({y}^{\prime \prime }\) & - & - & 0 & + \\
\cline{1-5}
\(y\) & 极大值 \(f\left( 0\right)  = 1\) & 减凹 ↘ & 拐点 \(\left( {\frac{1}{\sqrt{2}},\frac{1}{{\mathrm{e}}^{2}}}\right)\) & 减凸 ↘ \\
\cline{1-5}
\hline
\end{tabular}
}
\end{center}

渐近线 \(y = 0\) ;

(7)

\begin{center}
\adjustbox{max width=\textwidth}{
\begin{tabular}{|c|c|c|c|c|c|c|c|}
\hline
\phantom{X} & \(\left( {-\infty , - \frac{1}{5}}\right)\) & \(- \frac{1}{5}\) & \(\left( {-\frac{1}{5},0}\right)\) & 0 & \(\left( {0,\frac{2}{5}}\right)\) & \(\frac{2}{5}\) & \(\left( {\frac{2}{5}, + \infty }\right)\) \\
\cline{1-8}
\({y}^{\prime }\) & + & + & + & 不存在 & - & 0 & + \\
\cline{1-8}
\({y}^{\prime \prime }\) & - & 0 & + & 不存在 & + & + & + \\
\cline{1-8}
\(y\) & 增凹 ↗ & 拐点 \(\left( {-\frac{1}{5}, - \frac{6}{5}{\left( \frac{1}{5}\right) }^{\frac{2}{3}}}\right)\) & 增凸 ↗ & 极大值 \(f\left( 0\right)  = 0\) & 减凸 ↘ & 极小值 \(f\left( \frac{2}{5}\right)  =  - \frac{3}{5}{\left( \frac{2}{5}\right) }^{\frac{2}{3}}\) & 增凸 ↗ \\
\cline{1-8}
\hline
\end{tabular}
}
\end{center}

(8) 设 \({x}_{1} = \frac{1}{2} - \frac{3}{10}\sqrt{5},{x}_{2} = \frac{1}{2} + \frac{3}{10}\sqrt{5}\)

\begin{center}
\adjustbox{max width=\textwidth}{
\begin{tabular}{|c|c|c|c|c|c|c|}
\hline
\(x\) & \(\left( {-\infty ,{x}_{1}}\right)\) & \({x}_{1}\) & \(\left( {{x}_{1},0}\right)\) & 0 & \(\left( {0,\frac{1}{2}}\right)\) & \(\frac{1}{2}\) \\
\cline{1-7}
\({y}^{\prime }\) & - & - & - & 不存在 & + & 0 \\
\cline{1-7}
\({y}^{\prime \prime }\) & + & 0 & - & 不存在 & - & - \\
\cline{1-7}
\(y\) & 减凸 ↘ & 拐点 \(\left( {{x}_{1},f\left( {x}_{1}\right) }\right)\) & 减凹 ↘ & 极小值 \(f\left( 0\right)  = 0\) & 增凹 ↗ & 极大值 \(f\left( \frac{1}{2}\right)  = \frac{9}{4}{\left( \frac{1}{2}\right) }^{\frac{2}{3}}\) \\
\cline{1-7}
\(x\) & \multicolumn{2}{|c|}{\(\left( {\frac{1}{2},{x}_{2}}\right)\)} & \({x}_{2}\) & \(\left( {{x}_{2},2}\right)\) & 2 & \(\left( {2, + \infty }\right)\) \\
\cline{1-7}
\({y}^{\prime }\) & \multicolumn{2}{|c|}{-} & - & - & 0 & + \\
\cline{1-7}
\({y}^{\prime \prime }\) & \multicolumn{2}{|c|}{-} & 0 & + & + & + \\
\cline{1-7}
\(y\) & \multicolumn{2}{|c|}{减凹 ↘} & 拐点 \(\left( {{x}_{2},f\left( {x}_{2}\right) }\right)\) & 减凸 ↘ & 极小值 \(f\left( 2\right)  = 0\) & 增凸 ↗ \\
\cline{1-7}
\hline
\end{tabular}
}
\end{center}

\subsection{习题 6.7}

1. -1 . 20 .

2. 1.538 .

第六章总练习题

7. (1) e; (2) \(\frac{3}{2}\) ; (3) 0.

第七章 实数的完备性

习题 7.1

5. (1) 能; (2) (i) 不能, (ii) 能.

\subsection{习题 7.2}

1. (1) \(2,0;\left( 2\right) \frac{1}{2}, - \frac{1}{2};\left( 3\right)  + \infty , + \infty ;\left( 4\right) 2, - 2;\left( 5\right) \pi ,\pi ;\left( 6\right) 1,1\) .

第八章 不 定 积 分

\subsection{习题 8.1}

2. \(y = {x}^{2} + 1\) .

5. (1) \(x - \frac{{x}^{2}}{2} + \frac{{x}^{4}}{4} - 3\sqrt[3]{x} + C\) ; (2) \(\frac{{x}^{3}}{3} + \ln \left| x\right|  - \frac{4}{3}\sqrt{{x}^{3}} + C\) ;

(3) \(\sqrt{\frac{2x}{g}} + C\) ; (4) \(\frac{{4}^{x}}{\ln 4} + \frac{{9}^{x}}{\ln 9} + \frac{2 \cdot  {6}^{x}}{\ln 6} + C\) ;

(5) \(\frac{3}{2}\arcsin x + C\) ; (6) \(\frac{1}{3}\left( {x - \arctan x}\right)  + C\) ;

(7) \(\tan x - x + C\) ; (8) \(\frac{1}{4}\left( {{2x} - \sin {2x}}\right)  + C\) ;

(9) \(\sin x - \cos x + C\) ; (10) \(- \tan x - \cot x + C\) ;

(11) \(\frac{{90}^{t}}{\ln {90}} + C\) ; (12) \(\frac{8}{15}{x}^{\frac{15}{8}} + C\) ;

(13) 2 arcsin \(x + C\) ; (14) \(x - \frac{1}{2}\cos {2x} + C\) ;

( 15 ) \(\frac{1}{2}\left( {\sin x + \frac{1}{3}\sin {3x}}\right)  + C\) ; (16) \(\frac{1}{3}{\mathrm{e}}^{3x} - 3{\mathrm{e}}^{x} - 3{\mathrm{e}}^{-x} + \frac{1}{3}{\mathrm{e}}^{-{3x}} + C\) ;

(17) \(- \frac{2}{\ln 5} \cdot  {5}^{-x} + \frac{1}{5\ln 2} \cdot  {2}^{-x} + C\) ; (18) \(\ln \left| x\right|  - \frac{1}{4}{x}^{-4} + C\) .

6. (1) \(G\left( x\right)  + C\) ,其中 \(G\left( x\right)  = \left\{  \begin{array}{ll} 2 - {\mathrm{e}}^{-x}, & x \geq  0, \\  {\mathrm{e}}^{x}, & x < 0; \end{array}\right.\)

(2) \(G\left( x\right)  + C\) ,其中 \(G\left( x\right)  = \left\{  {\begin{array}{ll}  - \cos x + {4k}, & {2k\pi } \leq  x \leq  \left( {{2k} + 1}\right) \pi , \\  \cos x + {4k} + 2, & \left( {{2k} + 1}\right) \pi  \leq  x \leq  \left( {{2k} + 2}\right) \pi , \end{array}k \in  \mathbb{Z}}\right.\) .

7. \(\tan x - x + C\) .

\subsection{习题 8.2}

1. (1) \(\frac{1}{3}\sin \left( {{3x} + 4}\right)  + C\) ; (2) \(\frac{1}{4}{\mathrm{e}}^{2{x}^{2}} + C\) ;

(3) \(\frac{1}{2}\ln \left| {{2x} + 1}\right|  + C\) ; (4) \(\frac{{\left( 1 + x\right) }^{n + 1}}{n + 1} + C\) ;

(5) \(\arcsin \frac{x}{\sqrt{3}} + \frac{1}{\sqrt{3}}\arcsin \left( {\sqrt{3}x}\right)  + C\) ; (6) \(\frac{{2}^{{2x} + 2}}{\ln 2} + C\) ;

(7) \(- \frac{2}{9}\sqrt{{\left( 8 - 3x\right) }^{3}} + C\) ; (8) \(- \frac{3}{10}\sqrt[3]{{\left( 7 - 5x\right) }^{2}} + C\) ;

(9) \(- \frac{1}{2}\cos {x}^{2} + C\) ; (10) \(- \frac{1}{2}\cot \left( {{2x} + \frac{\pi }{4}}\right)  + C\) ;

(11) \(\tan \frac{x}{2} + C\) ; (12) \(\tan x - \sec x + C\) ;

(13) \(- \ln \left| {\csc x + \cot x}\right|  + C\) ; (14) \(- \sqrt{1 - {x}^{2}} + C\) ;

(15) \(\frac{1}{4}\arctan \frac{{x}^{2}}{2} + C\) ; (16) \(\ln \left| {\ln x}\right|  + C\) ;

(17) \(\frac{1}{10}{\left( 1 - {x}^{5}\right) }^{-2} + C\) ; (18) \(\frac{1}{8\sqrt{2}}\ln \left| \frac{{x}^{4} - \sqrt{2}}{{x}^{4} + \sqrt{2}}\right|  + C\) ;

(19) \(\ln \left| \frac{x}{1 + x}\right|  + C\) ; (20) \(\ln \left| {\sin x}\right|  + C\) ;

( 21 ) \(\sin x - \frac{2}{3}{\sin }^{3}x + \frac{1}{5}{\sin }^{5}x + C\) ; (22) \(\ln \left| {\tan x}\right|  + C\) ;

(23) arctan \({\mathrm{e}}^{x} + C\) ; (24) \(\ln \left| {{x}^{2} - {3x} + 8}\right|  + C\) ;

( 25 ) \(\ln \left| {x + 1}\right|  + \frac{2}{x + 1} - \frac{3}{2{\left( x + 1\right) }^{2}} + C\) ; (26) \(\ln \left| {x + \sqrt{{x}^{2} + {a}^{2}}}\right|  + C\) ;

(27) \(\frac{x}{{a}^{2}\sqrt{{x}^{2} + {a}^{2}}} + C\) ;

(28) \(- {\left( 1 - {x}^{2}\right) }^{\frac{1}{2}} + \frac{2}{3}{\left( 1 - {x}^{2}\right) }^{\frac{3}{2}} - \frac{1}{5}{\left( 1 - {x}^{2}\right) }^{\frac{5}{2}} + C\) ;

(29) \(- \frac{6}{7}{x}^{\frac{7}{6}} - \frac{6}{5}{x}^{\frac{5}{6}} - 2{x}^{\frac{1}{2}} - 6{x}^{\frac{1}{6}} - 3\ln \left| \frac{{x}^{\frac{1}{6}} - 1}{{x}^{\frac{1}{6}} + 1}\right|  + C\) ;

(30) \(x - 4\sqrt{x + 1} + 4\ln \left| {\sqrt{x + 1} + 1}\right|  + C\) ;

(31) \(- \frac{1}{4}\left( {\frac{1}{10100} + \frac{2}{101}x}\right) {\left( 1 - 2x\right) }^{100} + C\) ;

(32) \(\frac{1}{n}\ln \left| \frac{{x}^{n}}{1 + {x}^{n}}\right|  + C\) ; (33) \(\frac{1}{n}\left( {{x}^{n} - \ln \left| {1 + {x}^{n}}\right| }\right)  + C\) ;

(34) \(\ln \left| {\ln \ln x}\right|  + C\) ; (35) \(\ln x - \ln 2 \cdot  \ln \left| {\ln {4x}}\right|  + C\) ;

(36) \(\frac{\sqrt{{x}^{2} - 1}}{x} - \frac{{\left( {x}^{2} - 1\right) }^{\frac{3}{2}}}{3{x}^{3}} + C\) .

2. (1) \(x\arcsin x + \sqrt{1 - {x}^{2}} + C\) ; (2) \(x\ln x - x + C\) ;

(3) \({x}^{2}\sin x + {2x}\cos x - 2\sin x + C\) ;

(4) \(- \frac{1}{4{x}^{2}}\left( {2\ln x + 1}\right)  + C\) ; (5) \(x{\left( \ln x\right) }^{2} - {2x}\ln x + {2x} + C\) ;

(6) \(\frac{1}{2}\left( {{x}^{2} + 1}\right) \arctan x - \frac{x}{2} + C\) ; (7) \(x\ln \left( {\ln x}\right)  + C\) ;

(8) \(x{\left( \arcsin x\right) }^{2} + 2\sqrt{1 - {x}^{2}}\arcsin x - {2x} + C\) ;

(9) \(\frac{1}{2}\left( {\sec x\tan x + \ln \left| {\sec x + \tan x}\right| }\right)  + C\) ;

(10) \(\frac{1}{2}\left( {x\sqrt{{x}^{2} \pm  {a}^{2}} \pm  {a}^{2}\ln \left| {\sqrt{{x}^{2} \pm  {a}^{2}} + x}\right| }\right)  + C\) .

3. (1) \(\frac{1}{\alpha  + 1}{\left( f\left( x\right) \right) }^{\alpha  + 1} + C\) ; (2) \(\arctan \left( {f\left( x\right) }\right)  + C\) ;

(3) \(\ln \left| {f\left( x\right) }\right|  + C\) ; (4) \({\mathrm{e}}^{f\left( x\right) } + C\) .

5. (1) \(\frac{1}{2}{\tan }^{2}x + \ln \left| {\cos x}\right|  + C\) ; (2) \(\frac{1}{3}{\tan }^{3}x - \tan x + x + C\) ;

(3) \(\frac{x}{16} - \frac{1}{6}{\cos }^{3}x{\sin }^{3}x - \frac{1}{64}\sin {4x} + C\) .

6. (1) \({I}_{n} = \frac{1}{k}{x}^{n}{\mathrm{e}}^{kx} - \frac{n}{k}{I}_{n - 1}\) ; (2) \({I}_{n} = x{\left( \ln x\right) }^{n} - n{I}_{n - 1}\) ;

(3) \({I}_{n} = x{\left( \arcsin x\right) }^{n} + n\sqrt{1 - {x}^{2}}{\left( \arcsin x\right) }^{n - 1} - n\left( {n - 1}\right) {I}_{n - 2}\) ;

(4) \({I}_{n} = \frac{1}{{n}^{2} + {a}^{2}}\left\lbrack  {{\mathrm{e}}^{ax}{\sin }^{n - 1}x\left( {a\sin x - n\cos x}\right)  + n\left( {n - 1}\right) {I}_{n - 2}}\right\rbrack\) .

7. (1) \({\mathrm{e}}^{2x}\left( {\frac{1}{2}{x}^{3} - \frac{3}{4}{x}^{2} + \frac{3}{4}x - \frac{3}{8}}\right)  + C\) ;

(2) \(x\left\lbrack  {{\left( \ln x\right) }^{3} - 3{\left( \ln x\right) }^{2} + 6\ln x - 6}\right\rbrack   + C\) ;

(3) \(x{\left( \arcsin x\right) }^{3} + 3\sqrt{1 - {x}^{2}}{\left( \arcsin x\right) }^{2} - {6x}\arcsin x - 6\sqrt{1 - {x}^{2}} + C\) ;

(4) \(\frac{1}{10}{\mathrm{e}}^{x}\left( {{\sin }^{3}x - 3{\sin }^{2}x\cos x + 3\sin x - 3\cos x}\right)  + C\) .

\subsection{习题 8.3}

1. (1) \(\frac{{x}^{3}}{3} + \frac{{x}^{2}}{2} + x + \ln \left| {x - 1}\right|  + C\) ; (2) \(\ln \frac{{\left( x - 4\right) }^{2}}{\left| x - 3\right| } + C\) ;

(3) \(\frac{1}{6}\ln \frac{{\left( 1 + x\right) }^{2}}{{x}^{2} - x + 1} + \frac{1}{\sqrt{3}}\arctan \frac{{2x} - 1}{\sqrt{3}} + C\) ;

(4) \(\frac{\sqrt{2}}{8}\ln \left| \frac{{x}^{2} + \sqrt{2}x + 1}{{x}^{2} - \sqrt{2}x + 1}\right|  + \frac{\sqrt{2}}{4}\arctan \frac{\sqrt{2}x}{1 - {x}^{2}} + C\) ;

(5) \(\frac{1}{4}\ln \left| {x - 1}\right|  - \frac{1}{8}\ln \left( {{x}^{2} + 1}\right)  - \frac{1}{2}\arctan x - \frac{x - 1}{4\left( {{x}^{2} + 1}\right) } + C\) ;

(6) \(- \frac{{5x} + 3}{2\left( {2{x}^{2} + {2x} + 1}\right) } - \frac{5}{2}\arctan \left( {{2x} + 1}\right)  + C\) .

2. (1) \(\frac{1}{2}\arctan \left( {2\tan \frac{x}{2}}\right)  + C\) ; (2) \(\frac{\sqrt{6}}{6}\arctan \left( {\frac{\sqrt{6}}{2}\tan x}\right)  + C\) ;

(3) \(\frac{1}{2}\ln \left| {\cos x + \sin x}\right|  + \frac{x}{2} + C\) ;

(4) \(\frac{7}{8}\arcsin \frac{{2x} - 1}{\sqrt{5}} - \frac{{2x} + 3}{4}\sqrt{1 + x - {x}^{2}} + C\) ;

(5) \(\ln \left| {x + \frac{1}{2} + \sqrt{{x}^{2} + x}}\right|  + C\) ;

(6) \(\ln \left| \frac{1 + \sqrt{1 - {x}^{2}}}{x}\right|  - \frac{\sqrt{1 - {x}^{2}}}{x} + C\) .

\subsection{第八章总练习题}

1. (1) \(\frac{4}{5}{x}^{\frac{5}{4}} - \frac{24}{13}{x}^{\frac{13}{12}} - \frac{4}{3}{x}^{\frac{3}{4}} + C\) ;

(2) \(\frac{1}{2}{x}^{2}\arcsin x - \frac{1}{4}\arcsin x + \frac{1}{4}x\sqrt{1 - {x}^{2}} + C\) ;

(3) \(2\sqrt{x} - 2\ln \left( {1 + \sqrt{x}}\right)  + C\) ; (4) \(2{\mathrm{e}}^{\sin x}\left( {\sin x - 1}\right)  + C\) ;

(5) \(2{\mathrm{e}}^{\sqrt{x}}\left( {\sqrt{x} - 1}\right)  + C\) ; (6) arccos \(\frac{1}{x} + C\) ;

(7) \(\ln \left| {\cos x + \sin x}\right|  + C\) ; (8) \(\ln \left| {x - 2}\right|  - \frac{3}{x - 2} - \frac{1}{{\left( x - 2\right) }^{2}} + C\) ;

(9) \(\tan x + \frac{1}{3}{\tan }^{3}x + C\) ; (10) \(\frac{3}{8}x - \frac{1}{4}\sin {2x} + \frac{1}{32}\sin {4x} + C\) ;

( 11 ) \(\frac{2}{3}\ln \left| \frac{x - 2}{x + 1}\right|  + \frac{1}{x - 2} + C\) ;

(12) \(x\arctan \left( {1 + \sqrt{x}}\right)  - \sqrt{x} + \ln \left| {2 + x + 2\sqrt{x}}\right|  + C\) ;

(13) \(\frac{1}{4}{x}^{4} - \frac{1}{2}\ln \left( {{x}^{4} + 2}\right)  + C\) ;

( 15 ) \(\frac{1}{99}{\left( 1 - x\right) }^{-{99}} - \frac{1}{49}{\left( 1 - x\right) }^{-{98}} + \frac{1}{97}{\left( 1 - x\right) }^{-{97}} + C\) ;

( 16 ) \(- \frac{1}{x}\arcsin x - \ln \left| \frac{1 + \sqrt{1 - {x}^{2}}}{x}\right|  + C\) ; (17) \(\frac{{x}^{2} - 1}{2}\ln \left( \frac{1 + x}{1 - x}\right)  + x + C\) ;

( 18 ) \(2\sqrt{\tan x}\left( {1 + \frac{1}{5}{\tan }^{5}x}\right)  + C\) ; (19) \(\frac{{\mathrm{e}}^{x}}{1 + {x}^{2}} + C\) ;

(20) \({I}_{n} = \frac{2}{\left( {{2n} + 1}\right) {b}_{1}}\left\lbrack  {{v}^{n}\sqrt{u} + n\left( {{a}_{2}{b}_{1} - {a}_{1}{b}_{2}}\right) {I}_{n - 1}}\right\rbrack\) .

2. (1) \(\frac{1}{4}\ln \frac{{x}^{2} + x + 1}{{x}^{2} - x + 1} + \frac{1}{2\sqrt{3}}\arctan \frac{{2x} + 1}{\sqrt{3}} + \frac{1}{2\sqrt{3}}\arctan \frac{{2x} - 1}{\sqrt{3}} + C\) ;

(2) \(- \frac{{x}^{5} + 2}{{10}\left( {{x}^{10} + 2{x}^{5} + 2}\right) } - \frac{1}{10}\arctan \left( {{x}^{5} + 1}\right)  + C\) ;

(3) \(- \frac{{x}^{n}}{{2n}\left( {{x}^{2n} + 1}\right) } + \frac{1}{2n}\arctan {x}^{n} + C\) ;

(4) \(\frac{1}{2}x + \frac{1}{8}\cos {2x} + \frac{1}{8}\sin {2x} + \frac{1}{8}\ln \left( {\sin {2x} + 1}\right)  + C\) .

3. (1) \(\frac{12}{7}{\left( 1 + \sqrt[4]{x}\right) }^{\frac{7}{3}} - 3{\left( 1 + \sqrt[4]{x}\right) }^{\frac{4}{3}} + C\) ;

(2) \(\frac{1}{4}\ln \frac{\sqrt[4]{1 + {x}^{4}} + x}{\sqrt[4]{1 + {x}^{4}} - x} - \frac{1}{2}\arctan \frac{\sqrt[4]{1 + {x}^{4}}}{x} + C\) ;

(3) \(- \frac{3}{2\left( {{2x} + 2\sqrt{{x}^{2} - x + 1} - 1}\right) } - \frac{3}{2}\ln \left| {{2x} + 2\sqrt{{x}^{2} - x + 1} - 1}\right|  + 2\ln \left| {x + \sqrt{{x}^{2} - x + 1}}\right|  + C\) ;

(4) \(\frac{x}{\sqrt{1 - {x}^{4}}} + C\) .

5. (1) \({I}_{n} = \frac{\sin x}{\left( {n - 1}\right) {\cos }^{n - 1}x} + \frac{n - 2}{n - 1}{I}_{n - 2},\;n \geq  2\) ;

(2) \({I}_{n} = \frac{2}{n - 1}\sin \left( {n - 1}\right) x + {I}_{n - 2},\;n \geq  2\) .

\section{第九章 定 积 分}

\subsection{习题 9.1}

2. (1) \(\frac{1}{4};\left( 2\right) \mathrm{e} - 1;\left( 3\right) {\mathrm{e}}^{b} - {\mathrm{e}}^{a};\left( 4\right) \frac{1}{a} - \frac{1}{b}\) .

\subsection{习题 9.2}

1. (1) \(4;\;\left( 2\right) \frac{\pi }{2} - 1;\left( 3\right) \ln 2;\left( 4\right) \frac{\mathrm{e} + {\mathrm{e}}^{-1}}{2} - 1\) ;

(5) \(\sqrt{3} - \frac{\pi }{3}\) ; (6) \(\frac{44}{3}\) ; (7) \(4 - 2\ln 3\) ; (8) \(\frac{2}{3}\) .

2. (1) \(\frac{1}{4};\left( 2\right) \frac{1}{2};\left( 3\right) \frac{\pi }{4};\left( 4\right) \frac{2}{\pi }\) .

\subsection{习题 9.4}

6. a.

10. 提示: 证得存在第一个零点 \({x}_{1}\) 后,考察辅助函数 \(g\left( x\right)  = \left( {x - {x}_{1}}\right) f\left( x\right)\) .

11. 提示: \(f\) 凸,等价于曲线在任一切线的上方. (1) 取 \({x}_{0} = \frac{a + b}{2}\) ; (2) \(f\left( x\right)  \geq  f\left( t\right)  + {f}^{\prime }\left( t\right) \left( {x - t}\right)\) ,对 \(t\) 积分.

12. 提示: \(\frac{1}{k + 1} < \frac{1}{x} < \frac{1}{k}\) ,在 \(\left\lbrack  {k,k + 1}\right\rbrack\) 上积分.

\subsection{习题 9.5}

3. (1) 1 ; (2) 0 .

4. (1) \(\frac{2}{7};\left( 2\right) \frac{\pi }{3} + \frac{\sqrt{3}}{2};\left( 3\right) \frac{\pi {a}^{4}}{16};\left( 4\right) \frac{4}{3}\) ;

(5) arctan \(\mathrm{e} - \frac{\pi }{4}\) ; (6) \(\frac{\pi }{4}\) ; (7) \(\frac{\pi }{2} - 1\) ; (8) \(\frac{1}{2}\left( {{\mathrm{e}}^{\frac{\pi }{2}} + 1}\right)\) ;

(9) \(2 - \frac{2}{\mathrm{e}};\left( {10}\right) 2;\left( {11}\right) {a}^{3}\left( {\frac{\pi }{4} - \frac{2}{3}}\right) ;\left( {12}\right) \frac{\pi }{4}\) .

8. \(J\left( {{2m},{2n}}\right)  = \frac{\left( {{2n} - 1}\right) !!\left( {{2m} - 1}\right) !!}{{2}^{m + n}\left( {m + n}\right) !} \cdot  \frac{\pi }{2}\left( { = \frac{\pi \left( {2n}\right) !\left( {2m}\right) !}{{2}^{{2m} + {2n} + 1}m!n!\left( {m + n}\right) !}}\right)\) .

10. \(f\left( x\right)  - f\left( a\right) ;\;1 - \cos x\) .

12, 13. 提示: 使用积分第二中值定理.

16. 提示: \(F\left( x\right)  = {\int }_{a}^{x}f\left( t\right) \mathrm{d}t\) ,对 \({\int }_{a}^{b}f\left( x\right) g\left( x\right) \mathrm{d}x\) 进行分部积分.

\subsection{第九章总练习题}

1. 提示: \(f\) 凸, \(f\left( x\right)  \geq  f\left( {x}_{0}\right)  + {f}^{\prime }\left( {x}_{0}\right) \left( {x - {x}_{0}}\right) ,{x}_{0} = \frac{1}{a}{\int }_{0}^{a}\varphi \left( t\right) \mathrm{d}t,x = \varphi \left( t\right)\) ,并积分之.

3. 提示: \(\frac{1}{x}{\int }_{0}^{x}f\left( t\right) \mathrm{d}t = \frac{1}{x}{\int }_{0}^{\sqrt{x}}f\left( t\right) \mathrm{d}t + \frac{1}{x}{\int }_{\sqrt{x}}^{x}f\left( t\right) \mathrm{d}t\) ,并考察右边两项的极限.

4. 提示: \(x = {np} + {x}^{ * },0 < {x}^{ * } \leq  p\) ,利用周期函数的积分性质.

8. 提示: 与第 1 题类似,但需注意 \(\ln u\) 为凹函数.

9. 提示: 证明 \(\left\{  {a}_{n}\right\}\) 递减,有下界.

\section{第十章 定积分的应用}

\subsection{习题 10.1}

1. \(\frac{8}{3}\) . 2. \(\frac{1}{10}\left( {{99}\ln {10} - {81}}\right)\) . 3. \(\left( {{3\pi } + 2}\right) /\left( {{9\pi } - 2}\right)\) . 4. \(\frac{3}{8}\pi {a}^{2}\) . 5. \(\frac{3}{2}\pi {a}^{2}\) .

6. \(\frac{1}{4}\pi {a}^{2}\) . 7. \(\frac{1}{6}{ab}\) . 8. \(\frac{16}{35}\) . 9. \(\frac{5\pi }{24}\frac{\sqrt{3}}{4}\) . 10. 4ab \(\arcsin \frac{b}{\sqrt{{a}^{2} + {b}^{2}}}\left( {0 < b < a}\right)\) .

\subsection{习题 10.2}

1. \(\frac{400}{3}{\mathrm{\;{cm}}}^{3}\) .

2. (1) \(\frac{{\pi }^{2}}{2}\) ; (2) \(5{\pi }^{2}{a}^{3}\) ; (3) \(\frac{8}{3}\pi {a}^{3}\) ; (4) \(\frac{4}{3}\pi {a}^{2}b\) .

4. \(\frac{32}{105}\pi {a}^{3}.\;6.2{\pi }^{2}\)

\subsection{习题 10.3}

1. (1) \(\frac{8}{27}\left( {{10}\sqrt{10} - 1}\right)\) ; (2) \(1 + \frac{\sqrt{2}}{2}\ln \left( {\sqrt{2} + 1}\right)\) ; (3) \({6a}\) ;

(4) \(2{\pi }^{2}a\) ； (5) \(\frac{3}{2}{\pi a}\) ; (6) \({a\pi }\sqrt{1 + 4{\pi }^{2}} + \frac{a}{2}\ln \left( {{2\pi } + \sqrt{1 + 4{\pi }^{2}}}\right)\) .

2. (1) \(\frac{\sqrt{2}}{4}\) ; (3) \(\frac{\sqrt{2}}{4a}\) ; (4) \(\frac{2}{3a}\) .

3. \(a = 1,b = \sqrt{2}\) (或 \(a = \sqrt{2},b = 1\) ).

6. \(\frac{3}{4a},\frac{4a}{3},{\left( x - \frac{2}{3}a\right) }^{2} + {y}^{2} = \frac{16}{9}{a}^{2}\) .

8. \(\left( {-\ln \sqrt{2},\frac{\sqrt{2}}{2}}\right)\) .

\subsection{习题 10.4}

1. (1) \({2\pi }\left( {\sqrt{2} + \ln \left( {1 + \sqrt{2}}\right) }\right)\) ; (2) \(\frac{64}{3}\pi {a}^{2}\) ;

(3) \(a = b\) 时 \(S = {4\pi }{a}^{2},a < b\) 时 \(S = {2\pi a}\left( {a + \frac{{b}^{2}}{\sqrt{{b}^{2} - {a}^{2}}}\arcsin \frac{\sqrt{{b}^{2} - {a}^{2}}}{b}}\right)\) ,

\(a > b\) 时 \(S = {2\pi a}\left( {a + \frac{{b}^{2}}{\sqrt{{a}^{2} - {b}^{2}}}\ln \frac{\sqrt{{a}^{2} - {b}^{2}} + a}{b}}\right)\) ;

(4) \(4{\pi }^{2}{ar}\) .

2. \(S = {2\pi }{\int }_{\alpha }^{\beta }r\left( \theta \right) \sin \theta \sqrt{{r}^{2}\left( \theta \right)  + {r}^{\prime 2}\left( \theta \right) }\mathrm{d}\theta\) .

3. (1) \(\frac{32}{5}\pi {a}^{2}\) ; (2) \({4\pi }{a}^{2}\left( {2 - \sqrt{2}}\right)\) .

\subsection{习题 10.5}

1. \({14373.33}\mathrm{{kN}}\) .

3. \(- {1108.35}\mathrm{{kN}}\) . 4. \(\frac{kmM}{a\left( {a + l}\right) }\) . 5 . 5. \(\frac{k{M}^{2}}{{l}^{2}}\ln \frac{{\left( c + l\right) }^{2}}{c\left( {c + {2l}}\right) }\) .

6. \(\frac{2k\delta }{r}\) . 7. \({76969.02}\mathrm{\;{kJ}}\) . 8. \({3920}\mathrm{\;{kJ}}\) .

9. \(\frac{27}{7}k{a}^{\frac{7}{3}}{c}^{\frac{2}{3}}.\;{10}.\frac{4}{3}\pi {r}^{4}g\) .

\subsection{习题 10.6}

1. 0.6938, 0.6931 .

2.1.8569,1.8522,1.8519.

3. \({8.64}\left( {\mathrm{\;m}}^{2}\right)\) .

4. 矩形法平均:28.71 或 28.66 ；梯形法平均:28.68 ；抛物线法平均:28.67.

\section{第十一章 反常积分}

\subsection{习题 11.1}

1. (1) \(\frac{1}{2};\;\left( 2\right) 0;\;\left( 3\right) 2;\;\left( 4\right) 1 - \ln 2\) ;

(5) \(\frac{\pi }{4}\) ; (6) \(\frac{1}{2}\) ; (7) 发散；(8) 发散.

2. (1) \(p < 1\) 时收敛于 \(\frac{{\left( b - a\right) }^{1 - p}}{1 - p},p \geq  1\) 时发散;

(2)发散；(3)4；(4)1；

(5) -1; (6) \(\frac{\pi }{2}\) ; (7) \(\pi\) ; (8) 发散.

\subsection{习题 11.2}

4. (1)收敛；(2)收敛；(3)发散；

(4)收敛；(5) \(n > 1\) 时收敛, \(n \leq  1\) 时发散；

(6) \(n - m > 1\) 时收敛, \(n - m \leq  1\) 时发散.

5. (1) 条件收敛; (2)绝对收敛； (4)条件收敛.

\subsection{习题 11.3}

3. (1)发散；(2)收敛；(3)发散；(4)收敛；

(5)发散；(6) \(m < 3\) 时收敛, \(m \geq  3\) 时发散；

(7) \(0 < \alpha  < 1\) 时绝对收敛, \(1 \leq  \alpha  < 2\) 时条件收敛, \(\alpha  \geq  2\) 时发散; (8) 收敛.

4. (1) \({\left( -1\right) }^{n}n!;\;\left( 2\right) \frac{{2}^{{2n} + 1}{\left( n!\right) }^{2}}{\left( {{2n} + 1}\right) !}\) .

\subsection{第十一章总练习题}

3. (1) \(\frac{a}{{a}^{2} + {b}^{2}}\) ; (2) \(\frac{b}{{a}^{2} + {b}^{2}}\) ; (4) 0 .

4. \(0 < \lambda  \leq  1\) 时条件收敛; \(1 < \lambda  < 2\) 时绝对收敛; \(\lambda  \leq  0\) 或 \(\lambda  \geq  2\) 时发散.

6. (2) 提示: 证明充分性时问题归为证明 \(\mathop{\lim }\limits_{{x \rightarrow   + \infty }}{xf}\left( x\right)\) 存在,这可由 \({f}^{\prime }\left( x\right)\) 的定号性进而估计 \(\left| {{\int }_{u}^{+\infty }x{f}^{\prime }\left( x\right) \mathrm{d}x}\right| \left( {u\text{ 足够大 }}\right)\) 而得.

\section{上册索引}

垂直渐近线 62

存在域 9

半开半闭区间 4

D

待定系数法 177

保不等式性 47

单侧导数 85

被积表达式 162

比较原则 253,259

闭区间 4

闭区间套 150

变化率 84

变上限的定积分 205

\begin{center}
\adjustbox{max width=\textwidth}{
\begin{tabular}{|c|c|c|}
\hline
171 67 171 67 13 99 95 9 E 100 100 106 \(\mathbf{F}\) 22 248 13 13 59 8 8 88 171,209 9 68 10 5 11 G 100 & \phantom{X} & \phantom{X} \\
\cline{1-3}
递增数列 33 & \phantom{X} & 高阶微分107 \\
\cline{1-3}
第二换元公式 & \phantom{X} & 高阶无穷小量57 \\
\cline{1-3}
第二类间断点 & \phantom{X} & 根的存在定理72 \\
\cline{1-3}
第一换元公式 & \phantom{X} & 公式211 \\
\cline{1-3}
第一类间断点 & \phantom{X} & 拐点 \\
\cline{1-3}
定积分188 & \phantom{X} & 光滑曲线98 \\
\cline{1-3}
定义域9 & \phantom{X} & 归结原则50 \\
\cline{1-3}
对数函数 & \phantom{X} & \(\mathbf{H}\) \\
\cline{1-3}
对数螺线 & \phantom{X} & \phantom{X} \\
\cline{1-3}
对数求导法 & \phantom{X} & 海涅153 \\
\cline{1-3}
多值函数 & \phantom{X} & 函数值9 \\
\cline{1-3}
\phantom{X} & \phantom{X} & 弧长 229 弧微分233 \\
\cline{1-3}
二阶导数 & \phantom{X} & 换元积分法166 \\
\cline{1-3}
二阶可导 & \phantom{X} & J \\
\cline{1-3}
二阶微分 & \phantom{X} & \phantom{X} \\
\cline{1-3}
\phantom{X} & \phantom{X} & 积分变量162,188 \\
\cline{1-3}
\phantom{X} & \phantom{X} & 积分常数162 \\
\cline{1-3}
发散 248 & \phantom{X} & 中值定理207 \\
\cline{1-3}
发散数列 & \phantom{X} & 积分第一值定理202 \\
\cline{1-3}
反常积分 & \phantom{X} & 积分和188 \\
\cline{1-3}
反函数11 & \phantom{X} & 积分区间188 \\
\cline{1-3}
反三角函数 & \phantom{X} & 积分曲线163 \\
\cline{1-3}
非初等函数 & \phantom{X} & 积分型余项212 \\
\cline{1-3}
非正常极限 & \phantom{X} & 基本初等平面数 13 \\
\cline{1-3}
非正常上确界 & \phantom{X} & 基本周期17 \\
\cline{1-3}
非正常下确界 & \phantom{X} & 极大值88 \\
\cline{1-3}
费马定理 & \phantom{X} & 极限21,41,42 \\
\cline{1-3}
分部积分法 & \phantom{X} & 极小值88 \\
\cline{1-3}
分段函数 & \phantom{X} & 极值88 \\
\cline{1-3}
分段连续 & \phantom{X} & 间断点66 \\
\cline{1-3}
分割187 & \phantom{X} & 渐近线61 \\
\cline{1-3}
分划265 & \phantom{X} & 截面面积函数225 \\
\cline{1-3}
符号函数 & \phantom{X} & 介值性定理72 \\
\cline{1-3}
负无穷大 & \phantom{X} & 局部保号性47,69 \\
\cline{1-3}
复合函数 & \phantom{X} & 局部有界性47,69 \\
\cline{1-3}
\phantom{X} & \phantom{X} & 矩形法243 \\
\cline{1-3}
\phantom{X} & \phantom{X} & 聚点151 \\
\cline{1-3}
高阶导数 & \phantom{X} & 聚点定理152 \\
\cline{1-3}
\hline
\end{tabular}
}
\end{center}

\begin{center}
\adjustbox{max width=\textwidth}{
\begin{tabular}{|c|c|}
\hline
绝对收敛253 & \phantom{X} \\
\cline{1-2}
绝对值 & \(\mathbf{N}\) \\
\cline{1-2}
\(\mathbf{K}\) & 内函数 \\
\cline{1-2}
\phantom{X} & 逆映射12 \\
\cline{1-2}
开覆盖 153 & 牛顿一莱布尼茨 190 \\
\cline{1-2}
开区间 4 & 牛顿一莱布尼茨公式190 \\
\cline{1-2}
柯西判别法254 & 牛顿切线法145 \\
\cline{1-2}
柯西型余项 & \phantom{X} \\
\cline{1-2}
柯西中值定理 11 & 0 \\
\cline{1-2}
柯西准则 & 偶函数 17 \\
\cline{1-2}
可导 84 & \phantom{X} \\
\cline{1-2}
可导函数 86 & \(\mathbf{P}\) \\
\cline{1-2}
可积 188 & 抛物线法 2.244 \\
\cline{1-2}
可求长 229 & 迫敛性 \\
\cline{1-2}
可去间断点 & \phantom{X} \\
\cline{1-2}
可微 104 & Q \\
\cline{1-2}
可微函数 105 & 奇函数 17 \\
\cline{1-2}
空心 δ 邻域 5 & 切线方程 87 \\
\cline{1-2}
L & 区间 5 区间套 150 \\
\cline{1-2}
拉格朗日公式112 & 区间套定理 150 \\
\cline{1-2}
拉格朗日型余项129 & 曲边梯形 186 \\
\cline{1-2}
莱布尼茨公式101 & 曲率 233 \\
\cline{1-2}
黎曼 188 & 曲率半径 234 \\
\cline{1-2}
黎曼函数 & 曲率圆 234 \\
\cline{1-2}
黎曼和 188 & 曲率中心 234 \\
\cline{1-2}
黎曼积分188 & 全序 263 \\
\cline{1-2}
利普希茨 & 确界原理 6 \\
\cline{1-2}
连续 65 & \phantom{X} \\
\cline{1-2}
连续函数68 & S \\
\cline{1-2}
链式法则94 & 三角函数 13 \\
\cline{1-2}
邻域 5 & 三角形不等式 3 \\
\cline{1-2}
\(\mathbf{M}\) & 三歧性263 \\
\cline{1-2}
\phantom{X} & 194 \\
\cline{1-2}
麦克劳林公式127 & 217 \\
\cline{1-2}
密切圆 234 & 上极限156 \\
\cline{1-2}
幂函数 13 & 上界 \\
\cline{1-2}
闵可夫斯基221 & 上类 \\
\cline{1-2}
\phantom{X} & 上确界5 \\
\cline{1-2}
\hline
\end{tabular}
}
\end{center}

\begin{center}
\adjustbox{max width=\textwidth}{
\begin{tabular}{|c|c|c|c|}
\hline
上限188 & X & \phantom{X} & \phantom{X} \\
\cline{1-4}
施瓦兹221 & \phantom{X} & \phantom{X} & \phantom{X} \\
\cline{1-4}
收敛21 & 细度187 & \phantom{X} & \phantom{X} \\
\cline{1-4}
收敛准则36 & 瑕点250 & \phantom{X} & \phantom{X} \\
\cline{1-4}
数列21 & 瑕积分250 & \phantom{X} & \phantom{X} \\
\cline{1-4}
数轴 3 & 下和194 & \phantom{X} & \phantom{X} \\
\cline{1-4}
\phantom{X} & 下积分217 & \phantom{X} & \phantom{X} \\
\cline{1-4}
\(\mathbf{T}\) & 下极限156 & \phantom{X} & \phantom{X} \\
\cline{1-4}
泰勒定理129 & 下界5 & \phantom{X} & \phantom{X} \\
\cline{1-4}
泰勒公式129 & 下类265 & \phantom{X} & \phantom{X} \\
\cline{1-4}
梯形法243 & 下确界6 & \phantom{X} & \phantom{X} \\
\cline{1-4}
条件收敛253 & 下限188 & \phantom{X} & \phantom{X} \\
\cline{1-4}
跳跃间断点67 & 象 9 & \phantom{X} & \phantom{X} \\
\cline{1-4}
通项21 & 斜渐近线 61 & \phantom{X} & \phantom{X} \\
\cline{1-4}
同阶无穷大量60 & 心形线100 & \phantom{X} & \phantom{X} \\
\cline{1-4}
同阶无穷小量57 & Y & \phantom{X} & \phantom{X} \\
\cline{1-4}
凸函数 & \phantom{X} & \phantom{X} & \phantom{X} \\
\cline{1-4}
推广的确界原理8 & 严格凹函数 137 & \phantom{X} & \phantom{X} \\
\cline{1-4}
W & 严格单调函数 16 & \phantom{X} & \phantom{X} \\
\cline{1-4}
\phantom{X} & 严格(递)减函数16 & \phantom{X} & \phantom{X} \\
\cline{1-4}
外函数11 & 严格(递)增函数16 & \phantom{X} & \phantom{X} \\
\cline{1-4}
完备性150 & 严格凸函数 137 & \phantom{X} & \phantom{X} \\
\cline{1-4}
微分104 & 一阶微分形式的不变性 & 106 & \phantom{X} \\
\cline{1-4}
微积分学基本定理206 & 一致连续 74 & \phantom{X} & \phantom{X} \\
\cline{1-4}
106 & 一致连续性定理76 & \phantom{X} & \phantom{X} \\
\cline{1-4}
236 & 因变量 & \phantom{X} & \phantom{X} \\
\cline{1-4}
稳定点 88 & 映射 9 & \phantom{X} & \phantom{X} \\
\cline{1-4}
沃利斯 211 & 有界函数 15 & \phantom{X} & \phantom{X} \\
\cline{1-4}
无界集 5 & 有界集 5 & \phantom{X} & \phantom{X} \\
\cline{1-4}
无理数 1 & 有界数列 27 & \phantom{X} & \phantom{X} \\
\cline{1-4}
无穷大 4 & 有界性定理71 & \phantom{X} & \phantom{X} \\
\cline{1-4}
无穷大量 & 有理函数175 & \phantom{X} & \phantom{X} \\
\cline{1-4}
无穷积分248 & 有理数 1 & \phantom{X} & \phantom{X} \\
\cline{1-4}
无穷限反常积分 & \phantom{X} & \phantom{X} & \phantom{X} \\
\cline{1-4}
无穷小量56 & 有下界函数 15 & \phantom{X} & \phantom{X} \\
\cline{1-4}
无穷小数列25 & 有限覆盖定理153 & \phantom{X} & \phantom{X} \\
\cline{1-4}
无限开覆盖153 & 有限开覆盖153 & \phantom{X} & \phantom{X} \\
\cline{1-4}
无限区间 & 有限区间 4 & \phantom{X} & \phantom{X} \\
\cline{1-4}
\phantom{X} & 有限增量公式85 & \phantom{X} & \phantom{X} \\
\cline{1-4}
\hline
\end{tabular}
}
\end{center}

\begin{center}
\adjustbox{max width=\textwidth}{
\begin{tabular}{|c|c|}
\hline
右导数85 & 致密性定理36,153 \\
\cline{1-2}
右极限45 & 中间变量11 \\
\cline{1-2}
右连续68 & 中值定理111,112 \\
\cline{1-2}
右邻域5 & 周期函数17 \\
\cline{1-2}
域 263 & 驻点 88 \\
\cline{1-2}
原函数161 & 子列 31 \\
\cline{1-2}
原象9 & 自变量9 \\
\cline{1-2}
Z & 自变量增量 84 最大、最小值定理71 \\
\cline{1-2}
增量65 & 左导数85 \\
\cline{1-2}
振幅194 & 左极限45 \\
\cline{1-2}
正无穷大4 & 左连续68 \\
\cline{1-2}
值域 & 左邻域5 \\
\cline{1-2}
指数函数13 & \phantom{X} \\
\cline{1-2}
\hline
\end{tabular}
}
\end{center}

\section{第十二章 数项 级 数}

\subsection{习题 12.1}

1. (1) \(\frac{1}{5};\left( 2\right) \frac{3}{2};\left( 3\right) \frac{1}{4};\left( 4\right) 1 - \sqrt{2};\left( 5\right) 3\) .

6. (1) \(\frac{1}{a};\left( 2\right) 1;\left( 3\right) \frac{1}{2}\) .

7. (1) 收敛；(2) 发散；(3) 收敛；(4) 发散.

10. 提示: 对任意 \(n\) ,必有某 \(k\) ,使 \({n}_{k} < n \leq  {n}_{k + 1}\) ,则 \({S}_{{n}_{k}} < {S}_{n} \leq  {S}_{{n}_{k + 1}}\) 和 \({S}_{{n}_{k}} > {S}_{n} \geq  {S}_{{n}_{k + 1}}\) 必有其中之一成立.

\subsection{习题 12.2}

1. (1) 收敛；(2) 收敛；(3) 发散；(4) 收敛；(5) 收敛；

(6)发散；(7)发散；(8)收敛；(9)收敛；(10)收敛.

2. (1) 发散；(2) 发散；(3) 收敛；(4) 收敛；

(5)收敛；(6) \(a > b\) ,收敛； \(a < b\) ,发散.

9. (1) 收敛; (2) 发散; (3) 发散;

(4) \(p > 1\) ,收敛; \(p = 1,q > 1\) ,收敛; \(p = 1,q \leq  1\) ,发散; \(p < 1\) ,发散.

10. (1) 发散; (2) 收敛; (3) 收敛; (4) 收敛; (5) 发散; (6) 收敛(提示: \({3}^{\ln n} = {n}^{\ln 3}\) );

(7) 收敛.

12. (1) 收敛; (2) \(x > 1\) ,收敛; \(x \leq  1\) ,发散.

14. (1) 0 ; (2) 0 .

\subsection{习题 12.3}

1. (1) 绝对收敛; (2) 发散; (3) 当 \(p > 1\) 时绝对收敛,当 \(0 < p \leq  1\) 时条件收敛,当 \(p \leq  0\) 时发散;

(4)条件收敛；(5)发散；(6)条件收敛；(7)绝对收敛；(8) \(\left| x\right|  < \mathrm{e}\) 时绝对收敛, \(\left| x\right|  \geq  \mathrm{e}\) 时发散；

(9)条件收敛；(10)发散.

2. (1) 收敛；(2) 收敛；(3) 收敛.

5. (1) \(1 + 2{x}^{2} + \cdots  + n{x}^{{2n} - 2} + \cdots ;\left( 2\right) 1\) .

8. 提示: 先证 \(\sum {\left( -1\right) }^{m}\left( {\frac{1}{{m}^{2}} + \frac{1}{{m}^{2} + 1} + \cdots  + \frac{1}{{m}^{2} + {2m}}}\right)\) 收敛.

\subsection{第十二章总练习题}

7. 提示: 用 §2 习题 15 的结论.

\section{第十三章 函数列与函数项级数}

\subsection{习题 13.1}

1. (1)一致收敛；(2)一致收敛；(3)不一致收敛；(4)不一致收敛,内闭一致收敛；(5)不一致收敛, 内闭一致收敛.

3. (1) 一致收敛；(2) 一致收敛；(3) \(r > 1\) 时一致收敛, \(r = 1\) 时不一致收敛；(4) 一致收敛； (5)一致收敛；(6)不一致收敛.

9. (1) 一致收敛；(2) 不一致收敛；(3) 不一致收敛；(4) 一致收敛；(5) 一致收敛；(6) 不一致收敛.

\subsection{习题 13.2}

1. (1) \({f}_{n}\left( x\right)  \rightrightarrows  f\left( x\right)  = 1,{f}_{n}^{\prime }\left( x\right)  \rightrightarrows  g\left( x\right)  = 0\) ,三定理条件皆满足;

(2) \({f}_{n}\left( x\right)  \rightrightarrows  f\left( x\right)  = x,\left\{  {f}_{n}^{\prime }\right\}\) 不一致收敛,定理 13.9 和 13.10 条件满足;

(3) \(\left\{  {f}_{n}\right\}\) 与 \(\left\{  {f}_{n}^{\prime }\right\}\) 都不一致收敛,三定理条件均不满足,但定理 13.9 结论仍成立.

4. \(\mathop{\sum }\limits_{{n = 1}}^{\infty }\frac{{x}^{n}}{{n}^{3}}\) .

5. \(\mathop{\sum }\limits_{{n = 1}}^{\infty }\frac{\sin {nx}}{{n}^{2}\sqrt{n}}\) .

6. \(\frac{1}{2}\) .

9. (1) \({f}_{n} \rightrightarrows  f\left( x\right)  = 0\) ; (2)(i) \(\left\{  {f}_{n}\right\}\) 不一致收敛,极限函数不连续、不可微但可积; (ii) \({f}_{n} \rightrightarrows  f\left( x\right)  = 1\) .

\subsection{第十三章总练习题}

1. (1) \(k < 1\) 时一致收敛; (2) \(k < 1\) 时一致收敛.

\section{第十四章 幂 级 数}

\subsection{习题 14.1}

1. (1) \(R = 1,\left( {-1,1}\right) ;\left( 2\right) R = 2,\left\lbrack  {-2,2}\right\rbrack  ;\left( 3\right) R = 4,\left( {-4,4}\right) ;\left( 4\right) R =  + \infty ,\left( {-\infty , + \infty }\right) ;\left( 5\right) R =\)  \(+ \infty ,\left( {-\infty , + \infty }\right)\) ; (6) \(R = \frac{1}{3},\left\lbrack  {-\frac{4}{3}, - \frac{2}{3}}\right)\) ; (7) \(R = 1,\left( {-1,1}\right)\) ; (8) \(R = 1,\left\lbrack  {-1,1}\right\rbrack\) .

2. (1) \(\frac{1}{2}\ln \frac{1 + x}{1 - x},x \in  \left( {-1,1}\right)\) ; (2) \(\frac{2x}{{\left( 1 - x\right) }^{3}},x \in  \left( {-1,1}\right)\) ; (3) \(\frac{x + {x}^{2}}{{\left( 1 - x\right) }^{3}},x \in  \left( {-1,1}\right)\) .

8. (1) \(\left( {-R,R}\right) ,R = \max \{ a,b\} ;\left( 2\right) \left( {-\frac{1}{\mathrm{e}},\frac{1}{\mathrm{e}}}\right)\) .

9. (1) \(R = \frac{1}{4};\left( 2\right) R = 1\) .

10. (1) \(S\left( x\right)  = \left\{  \begin{array}{ll} \frac{1 - x}{x}\ln \left( {1 - x}\right)  + 1, & x \in  \lbrack  - 1,0) \cup  \left( {0,1}\right) , \\  0, & x = 0, \\  1, & x = 1; \end{array}\right.\)

(2) \(S\left( x\right)  =  - \frac{{\left( 1 - x\right) }^{2}}{2{x}^{2}}\ln \left( {1 - x}\right)  + \frac{3}{4} - \frac{1}{2x},x \in  \lbrack  - 1,0) \cup  \left( {0,1}\right) ;x = 1\) 时 \(S\left( x\right)  = \frac{1}{4};x = 0\) 时 \(S\left( x\right)  = 0\) ,

(3) \(S\left( x\right)  = \left\{  \begin{array}{ll} \frac{1}{{\left( 1 - x\right) }^{2}} - \frac{4}{\left( 1 - x\right) } - \frac{4}{x}\ln \left( {1 - x}\right) , & x \in  \left( {-1,0}\right)  \cup  \left( {0,1}\right) , \\  1, & x = 0. \end{array}\right.\)

11. (1) \(R = 1\) ; (2) \(2{a}_{0} + {2d}\) ( \(d\) 为公差).

\subsection{习题 14.2}

2. (1) \(\mathop{\sum }\limits_{{n = 0}}^{\infty }\frac{{x}^{2n}}{n!},x \in  \left( {-\infty ,\infty }\right)\) ; (2) \(\mathop{\sum }\limits_{{n = 0}}^{\infty }{x}^{n + {10}},x \in  \left( {-1,1}\right)\) ;

(3) \(x + \mathop{\sum }\limits_{{n = 1}}^{\infty }\frac{\left( {{2n} - 1}\right) !!}{n!}{x}^{n + 1},x \in  \left\lbrack  {-\frac{1}{2},\frac{1}{2}}\right)\) ;

(4) \(\mathop{\sum }\limits_{{n = 1}}^{\infty }{\left( -1\right) }^{n + 1}\frac{{2}^{{2n} - 1}{x}^{2n}}{\left( {2n}\right) !},x \in  \left( {-\infty , + \infty }\right)\) ;

(5) \(\mathop{\sum }\limits_{{n = 0}}^{\infty }\left( {\mathop{\sum }\limits_{{k = 0}}^{n}\frac{1}{k!}}\right) {x}^{n},x \in  \left( {-1,1}\right)\) ;

(6) \(\frac{1}{3}\mathop{\sum }\limits_{{n = 0}}^{\infty }\left( {1 - {\left( -1\right) }^{n}{2}^{n}}\right) {x}^{n},x \in  \left( {-\frac{1}{2},\frac{1}{2}}\right)\) ;

(7) \(\mathop{\sum }\limits_{{n = 0}}^{\infty }\frac{{\left( -1\right) }^{n}}{\left( {{2n} + 1}\right) !}\frac{{x}^{{2n} + 1}}{\left( 2n + 1\right) },x \in  \left( {-\infty , + \infty }\right)\) ;

(8) \(\mathop{\sum }\limits_{{n = 0}}^{\infty }\frac{{\left( -1\right) }^{n}\left( {1 - n}\right) }{n!}{x}^{n},x \in  \left( {-\infty , + \infty }\right)\) ;

(9) \(x + \mathop{\sum }\limits_{{n = 1}}^{\infty }\frac{{\left( -1\right) }^{n}}{\left( {2n}\right) !!}\frac{\left( {{2n} - 1}\right) !!}{\left( 2n + 1\right) }{x}^{{2n} + 1},x \in  \left\lbrack  {-1,1}\right\rbrack\) .

3. (1) \(8 + {15}\left( {x - 1}\right)  + {17}{\left( x - 1\right) }^{2} + 7{\left( x - 1\right) }^{3};\left( 2\right) \mathop{\sum }\limits_{{n = 0}}^{\infty }{\left( -1\right) }^{n}{\left( x - 1\right) }^{n}\) ;

(3) \(1 + \frac{3}{2}\left( {x - 1}\right)  + \mathop{\sum }\limits_{{n = 0}}^{\infty }{\left( -1\right) }^{n}\frac{\left( {2n}\right) !}{{\left( n!\right) }^{2}}\frac{3}{\left( {n + 1}\right) \left( {n + 2}\right) {2}^{n}}{\left( \frac{x - 1}{2}\right) }^{n + 2},x \in  \left\lbrack  {0,2}\right\rbrack\) .

4. (1) \(\frac{1}{2}\mathop{\sum }\limits_{{n = 0}}^{\infty }\left( {n + 1 - \frac{1 + {\left( -1\right) }^{n}}{2}}\right) {x}^{n},x \in  \left( {-1,1}\right)\) ;

(2) \(\mathop{\sum }\limits_{{n = 0}}^{\infty }{\left( -1\right) }^{n}\frac{{x}^{{2n} + 2}}{\left( {{2n} + 1}\right) \left( {{2n} + 2}\right) },x \in  \left\lbrack  {-1,1}\right\rbrack\) .

5. \(\mathop{\sum }\limits_{{n = 1}}^{\infty }\frac{2}{{2n} - 1}{\left( \frac{x - 1}{x + 1}\right) }^{{2n} - 1},x \in  \left( {0, + \infty }\right)\) .

\subsection{第十四章总练习题}

2. (1) \(x + \mathop{\sum }\limits_{{n = 2}}^{\infty }{\left( -1\right) }^{n}\frac{1}{n\left( {n - 1}\right) }{x}^{n},x \in  ( - 1,1\rbrack\) ; (2) \(\frac{1}{4}\mathop{\sum }\limits_{{n = 2}}^{\infty }{\left( -1\right) }^{n}\frac{{3}^{{2n} - 1} - 3}{\left( {{2n} - 1}\right) !}{x}^{{2n} - 1}\) , \(x \in  \left( {-\infty , + \infty }\right) ;\left( 3\right) \mathop{\sum }\limits_{{n = 0}}^{\infty }\frac{{\left( -1\right) }^{n}}{\left( {2n}\right) !}\frac{{x}^{{4n} + 1}}{{4n} + 1},x \in  \left( {-\infty , + \infty }\right) .\)

3. (1) \(\frac{1 + x}{{\left( 1 - x\right) }^{3}},x \in  \left( {-1,1}\right) ;\left( 2\right) \frac{2 + {x}^{2}}{{\left( 2 - {x}^{2}\right) }^{2}},x \in  \left( {-\sqrt{2},\sqrt{2}}\right)\) ;

(3) \(\frac{1}{{\left( 2 - x\right) }^{2}},x \in  \left( {0,2}\right)\) ; (4) \(\frac{{x}^{2}}{2}\arctan x + \frac{1}{2}\arctan x - \frac{x}{2},x \in  \left\lbrack  {-1,1}\right\rbrack\) .

4. (1) \(1;\left( 2\right) \frac{1}{3}\ln 2 + \frac{\pi }{3\sqrt{3}}\) .

6. (1) \(\frac{1}{2};\left( 2\right)  - \frac{1}{6}\) .

\section{第十五章 傅里叶级数}

\subsection{习题 15.1}

1. (1) (i) \(2\mathop{\sum }\limits_{{n = 1}}^{\infty }\frac{{\left( -1\right) }^{n + 1}}{n}\sin {nx},\;\) (ii) \(\pi  - 2\mathop{\sum }\limits_{{n = 1}}^{\infty }\frac{\sin {nx}}{n}\) ;

(2) (i) \(\frac{{\pi }^{2}}{3} + 4\mathop{\sum }\limits_{{n = 1}}^{\infty }{\left( -1\right) }^{n}\frac{1}{{n}^{2}}\cos {nx}\) ,(ii) \(\frac{4}{3}{\pi }^{2} + 4\mathop{\sum }\limits_{{n = 1}}^{\infty }\left( {\frac{\cos {nx}}{{n}^{2}} - \frac{\pi \sin {nx}}{n}}\right)\) ;

(3) \(\frac{b - a}{4}\pi  + \frac{2\left( {a - b}\right) }{\pi }\mathop{\sum }\limits_{{n = 1}}^{\infty }\frac{\cos \left( {{2n} - 1}\right) x}{{\left( 2n - 1\right) }^{2}} + \left( {a + b}\right) \mathop{\sum }\limits_{{n = 1}}^{\infty }{\left( -1\right) }^{n + 1}\frac{\sin {nx}}{n}\) .

3. \(\mathop{\sum }\limits_{{n = 1}}^{\infty }\frac{1}{{2n} - 1}\sin \left( {{2n} - 1}\right) x\) .

7. (1) \(\mathop{\sum }\limits_{{n = 1}}^{\infty }\frac{\sin {nx}}{n},x \in  \left( {0,{2\pi }}\right) ;\left( 2\right) \frac{2\sqrt{2}}{\pi }\left( {1 - 2\mathop{\sum }\limits_{{n = 1}}^{\infty }\frac{\cos {nx}}{4{n}^{2} - 1}}\right) ,x \in  \left\lbrack  {-\pi ,\pi }\right\rbrack\) ;

(3) (i) \(\frac{4}{3}a{\pi }^{2} + {b\pi } + c + \mathop{\sum }\limits_{{n = 1}}^{\infty }\frac{4a}{{n}^{2}}\cos {nx} - \frac{{4\pi a} + {2b}}{n}\sin {nx}\) ,

(ii) \(\frac{a}{3}{\pi }^{2} + c + \mathop{\sum }\limits_{{n = 1}}^{\infty }\left( {\frac{{\left( -1\right) }^{n}{4a}}{{n}^{2}}\cos {nx} - \frac{{\left( -1\right) }^{n}{2b}}{n}\sin {nx}}\right)\) ;

(4) \(\frac{\operatorname{sh}\pi }{\pi } + \frac{2}{\pi }\mathop{\sum }\limits_{{n = 0}}^{\infty }\frac{\operatorname{sh}\pi }{1 + {n}^{2}}{\left( -1\right) }^{n}\cos {nx};\;\left( 5\right) \frac{2}{\pi }\operatorname{sh}\pi \mathop{\sum }\limits_{{n = 1}}^{\infty }\frac{{\left( -1\right) }^{n - 1}n}{{n}^{2} + 1}\sin {nx}\) .

8. \(\mathop{\sum }\limits_{{n = 1}}^{\infty }\frac{\cos {nx}}{{n}^{2}}\) .

\subsection{习题 15.2}

1. (1) \(\frac{2}{\pi } + \frac{4}{\pi }\mathop{\sum }\limits_{{n = 1}}^{\infty }\frac{{\left( -1\right) }^{n + 1}}{4{n}^{2} - 1}\cos {2nx}\) ; (2) \(\frac{1}{2} - \frac{1}{\pi }\mathop{\sum }\limits_{{n = 1}}^{\infty }\frac{\sin {2\pi nx}}{n}\) ;

(3) \(\frac{3}{8} - \frac{1}{2}\cos {2x} + \frac{1}{8}\cos {4x}\) ; (4) \(\frac{4}{\pi }\mathop{\sum }\limits_{{n = 0}}^{\infty }\left\lbrack  {{\left( -1\right) }^{n}\frac{\cos \left( {{2n} + 1}\right) x}{{2n} + 1}}\right\rbrack\) .

2. \(\frac{2}{3} + \frac{3}{{\pi }^{2}}\mathop{\sum }\limits_{{n = 1}}^{\infty }\left( {{\left( -1\right) }^{n}\cos \frac{n\pi }{3} - 1}\right) \frac{1}{{n}^{2}}\cos \frac{2n\pi x}{3}\) .

3. \(\frac{4}{\pi }\mathop{\sum }\limits_{{n = 1}}^{\infty }\frac{\cos \left( {{2n} - 1}\right) x}{{\left( 2n - 1\right) }^{2}}\) .

4. \(\frac{8}{\pi }\mathop{\sum }\limits_{{n = 1}}^{\infty }\frac{n}{4{n}^{2} - 1}\sin {nx}\) .

5. \(\frac{8}{{\pi }^{2}}\mathop{\sum }\limits_{{n = 0}}^{\infty }\frac{1}{{\left( 2n + 1\right) }^{2}}\cos \frac{\left( {{2n} + 1}\right) {\pi x}}{2}\) .

6. \(\frac{1}{3} + \frac{4}{{\pi }^{2}}\mathop{\sum }\limits_{{n = 1}}^{\infty }\frac{\cos {n\pi x}}{{n}^{2}}\) .

7. (1) \(\frac{4}{\pi }\mathop{\sum }\limits_{{n = 1}}^{\infty }\frac{{\left( -1\right) }^{n + 1}}{{\left( 2n - 1\right) }^{2}}\sin \left( {{2n} - 1}\right) x\) ; (2) \(\frac{4}{\pi }\mathop{\sum }\limits_{{n = 1}}^{\infty }\frac{\cos \left( {{2n} - 1}\right) x}{{\left( 2n - 1\right) }^{2}}\) .

\subsection{第十五章总练习题}

4. (1) \({\alpha }_{n} = {a}_{n},{\beta }_{n} =  - {b}_{n};\left( 2\right) {\alpha }_{n} =  - {a}_{n},{\beta }_{n} = {b}_{n}\) .

\section{第十六章 多元函数的极限与连续}

习题 16.1

7. (1) \(\frac{9}{16};\left( 2\right) \frac{2xy}{{x}^{2} + {y}^{2}};\left( 3\right) {t}^{2}\left( {{x}^{2} + {y}^{2} - {xy}\tan \frac{x}{y}}\right)\) .

9. (1) \(y \neq   \pm  x\) ; (2) \(\left( {x,y}\right)  \neq  \left( {0,0}\right)\) ; (3) \({xy} \geq  0\) ; (4) \(\{ \left( {x,y}\right) \left| \right| x\left| { \leq  1,}\right| y \mid   \geq  1\}\) ;

(5) \(\left\{  {\left( {x,y}\right)  \mid  x > 0,y > 0}\right\}\) ; (6) \(\left\{  {\left( {x,y}\right)  \mid  {2n\pi } \leq  {x}^{2} + {y}^{2} \leq  \left( {{2n} + 1}\right) \pi ,n = 0,1,2,\cdots }\right\}\) ;

(7) \(\{ \left( {x,y}\right)  \mid  y > x\}\) ; (8) 全平面; (9) 整个三维空间; (10) \(\left\{  \left( {x,y,z}\right) \right\}  \left\{  {{r}^{2} < {x}^{2} + {y}^{2} + {z}^{2} \leq  {R}^{2}}\right\}\) .

\subsection{习题 16.2}

1. (1) 0; (2) + ∞; (3) 2; (4) + ∞; (5) ∞; (6) 0; (7) 1.

2. (1) 重极限不存在, \(\mathop{\lim }\limits_{{x \rightarrow  0}}\mathop{\lim }\limits_{{y \rightarrow  0}}f\left( {x,y}\right)  = 0,\mathop{\lim }\limits_{{y \rightarrow  0}}\mathop{\lim }\limits_{{x \rightarrow  0}}f\left( {x,y}\right)  = 1\) ;

(2) \(\mathop{\lim }\limits_{{\left( {x,y}\right)  \rightarrow  \left( {0,0}\right) }}f\left( {x,y}\right)  = 0\) ,两个累次极限均不存在；

(3) 重极限不存在, \(\mathop{\lim }\limits_{{x \rightarrow  0}}\mathop{\lim }\limits_{{y \rightarrow  0}}f\left( {x,y}\right)  = \mathop{\lim }\limits_{{y \rightarrow  0}}\mathop{\lim }\limits_{{x \rightarrow  0}}f\left( {x,y}\right)  = 0\) ;

(4)重极限不存在, \(\mathop{\lim }\limits_{{x \rightarrow  0}}\mathop{\lim }\limits_{{y \rightarrow  0}}f\left( {x,y}\right)  = \mathop{\lim }\limits_{{y \rightarrow  0}}\mathop{\lim }\limits_{{x \rightarrow  0}}f\left( {x,y}\right)  = 0\) ；

(5) \(\mathop{\lim }\limits_{{\left( {x,y}\right)  \rightarrow  \left( {0,0}\right) }}f\left( {x,y}\right)  = 0,\mathop{\lim }\limits_{{x \rightarrow  0}}\mathop{\lim }\limits_{{y \rightarrow  0}}f\left( {x,y}\right)  = 0\) ,另一累次极限不存在;

(6) 重极限不存在, \(\mathop{\liminf }\limits_{{x \rightarrow  0}}f\left( {x,y}\right)  = \mathop{\liminf }\limits_{{y \rightarrow  0}}\left( {x,y}\right)  = 0\) ;

(7)重极限与累次极限均不存在.

7. (1) 0; (2) 0; (3) 1; (4) e.

\subsection{习题 16.3}

1. (1) 间断曲线为圆族 \({x}^{2} + {y}^{2} = \left( {{2n} + 1}\right) \frac{\pi }{2},n = 0,1,2,\cdots\) ;

(2)间断曲线为直线族 \(x + y = n,n = 0, \pm  1, \pm  2,\cdots\) ;

(3) 不连续点集合 \(\{ \left( {x,y}\right)  \mid  x \neq  0,y = 0\}\) ; (4) 在 \({\mathbf{R}}^{2}\) 上连续;

(5) 仅在直线 \(y = 0\) 上连续; *(6) 在 \({\mathbf{R}}^{2}\) 上连续;

(7) 在定义域上连续; *(8) 在定义域上连续.

\subsection{第十六章总练习题}

3. (1) 存在; (2) 不存在.

\section{第十七章 多元函数微分学}

\subsection{习题 17.1}

1. (1) \({z}_{x} = {2xy},{z}_{y} = {x}^{2}\) ; (2) \({z}_{x} =  - y\sin x,{z}_{y} = \cos x\) ;

(3) \({z}_{x} = \frac{-x}{{\left( {x}^{2} + {y}^{2}\right) }^{3/2}},{z}_{y} = \frac{-y}{{\left( {x}^{2} + {y}^{2}\right) }^{3/2}}\) ; (4) \({z}_{x} = \frac{1}{x + {y}^{2}},{z}_{y} = \frac{2y}{x + {y}^{2}}\) ;

(5) \({z}_{x} = y{\mathrm{e}}^{xy},{z}_{y} = x{\mathrm{e}}^{xy};\left( 6\right) {z}_{x} =  - \frac{y}{{x}^{2} + {y}^{2}},{z}_{y} = \frac{x}{{x}^{2} + {y}^{2}}\) ;

(7) \({z}_{x} = y\left( {1 + {xy}\cos \left( {xy}\right) }\right)  \cdot  {\mathrm{e}}^{\sin \left( {xy}\right) },{z}_{y} = x\left( {1 + {xy}\cos \left( {xy}\right) }\right) {\mathrm{e}}^{\sin \left( {xy}\right) }\) ;

(8) \({u}_{x} =  - \frac{y}{{x}^{2}} - \frac{1}{z},{u}_{y} = \frac{1}{x} - \frac{z}{{y}^{2}},{u}_{z} = \frac{1}{y} + \frac{x}{{z}^{2}}\) ;

(9) \({u}_{x} = {yz}{\left( xy\right) }^{z - 1},{u}_{y} = {xz}{\left( xy\right) }^{z - 1},{u}_{z} = {\left( xy\right) }^{z}\ln \left( {xy}\right)\) ;

(10) \({u}_{x} = {y}^{z}{x}^{{y}^{z} - 1},{u}_{y} = z{y}^{z - 1}{x}^{{y}^{z}}\ln x,{u}_{z} = {y}^{z}{x}^{{y}^{z}}\ln x\ln y\) .

2. \({f}_{x}\left( {x,1}\right)  = 1\) .

3. \({f}_{x}\left( {0,0}\right)  = 0,{f}_{y}\left( {0,0}\right)\) 不存在.

8. (1) \({\left. \mathrm{d}z\right| }_{\left( 0,0\right) } = 0,{\left. \mathrm{d}z\right| }_{\left( 1,1\right) } =  - 4\mathrm{\;d}x - 4\mathrm{\;d}y;\left( 2\right) {\left. \mathrm{d}z\right| }_{\left( 1,0\right) } = 0,{\left. \mathrm{d}z\right| }_{\left( 0,1\right) } = \mathrm{d}x\) .

9. (1) \(\mathrm{d}z = y\cos \left( {x + y}\right) \mathrm{d}x + \left( {\sin \left( {x + y}\right)  + y\cos \left( {x + y}\right) }\right) \mathrm{d}y\) ;

(2) \(\mathrm{d}u = {\mathrm{e}}^{yz}\mathrm{\;d}x + \left( {{xz}{\mathrm{e}}^{yz} + 1}\right) \mathrm{d}y + \left( {{xy}{\mathrm{e}}^{yz} - {\mathrm{e}}^{-z}}\right) \mathrm{d}z\) .

10. \(x - y + {2z} = \frac{\pi }{2},2\left( {1 - x}\right)  = 2\left( {y - 1}\right)  = \frac{\pi }{4} - z\) .

11. \({9x} + y - z - {27} = 0,x - 3 = 9\left( {y - 1}\right)  = 9\left( {1 - z}\right)\) .

12. \(\left( {-3, - 1,3}\right) ,x + {3y} + z + 3 = 0,3\left( {x + 3}\right)  = y + 1 = 3\left( {z - 3}\right)\) .

13. (1) 108. 972; (2) 0.5023 .

14. \({2576}{\mathrm{\;{cm}}}^{3}\) .

18. \(\frac{\pi }{4}\) .

20. \({0.26}\% ,{0.02}\) .

\subsection{习题 17.2}

1. (1) \(\frac{\mathrm{d}z}{\mathrm{\;d}x} = \frac{\left( {1 + x}\right) {\mathrm{e}}^{x}}{1 + {x}^{2}{\mathrm{e}}^{2x}}\) ; (2) \({z}_{x} = \left( {1 + \frac{{x}^{2} + {y}^{2}}{xy}}\right) \frac{{x}^{2} - {y}^{2}}{{x}^{2}y}{\mathrm{e}}^{\frac{{x}^{2} + {y}^{2}}{xy}},{z}_{y} = \left( {1 + \frac{{x}^{2} + {y}^{2}}{xy}}\right) \frac{{y}^{2} - {x}^{2}}{x{y}^{2}}{\mathrm{e}}^{\frac{{x}^{2} + {y}^{2}}{xy}}\) ;

(3) \(\frac{\mathrm{d}z}{\mathrm{\;d}t} = 4{t}^{3} + 3{t}^{2} + {2t}\) ; (4) \({z}_{u} = \frac{u}{{v}^{2}}\left( {2\ln \left( {{3u} - {2v}}\right)  + \frac{3u}{{3u} - {2v}}}\right) ,{z}_{v} =  - \frac{2{u}^{2}}{{v}^{2}}\left( {\frac{1}{v}\ln \left( {{3u} - {2v}}\right)  + \frac{1}{{3u} - {2v}}}\right)\) ;

(5) \({u}_{x} = {f}_{1} + y{f}_{2},{u}_{y} = {f}_{1} + x{f}_{2}\) ; (6) \({u}_{x} = \frac{1}{y}{f}_{1},{u}_{y} =  - \frac{x}{{y}^{2}}{f}_{1} + \frac{1}{z}{f}_{2},{u}_{z} =  - \frac{y}{{z}^{2}}{f}_{2}\) .

2. \(\mathrm{d}z = {\left( x + y\right) }^{xy}\left\{  {\left\lbrack  {\frac{xy}{x + y} + y\ln \left( {x + y}\right) }\right\rbrack  \mathrm{d}x + \left\lbrack  {\frac{xy}{x + y} + x\ln \left( {x + y}\right) }\right\rbrack  \mathrm{d}y}\right\}\) .

6. \({F}_{x}\left( {0,0}\right)  = 4{f}^{\prime }\left( 0\right) ;{F}_{t}\left( {0,0}\right)  = 0\) .

习题 17.3

1. 5 .

2. \(\frac{98}{13}\) .

3. \(\left( {-4,2, - 4}\right) ,6;\left( {3, - 5,0}\right) ,\sqrt{34}\) .

4. \(- \frac{1}{{r}^{2}}\left( {x - a,y - b,z - c}\right) ,r = 1\) .

5. \(- 2\left( {\frac{1}{a},\frac{1}{b},\frac{1}{c}}\right)\) .

7. (1) \(\frac{1}{r}\left( {x,y,z}\right) ;\left( 2\right)  - \frac{1}{{r}^{3}}\left( {x,y,z}\right)\) .

8. (1) \({z}^{2} = {xy};\left( 2\right) {x}^{2} = {yz},{y}^{2} = {zx};\left( 3\right) x = y = z\) .

习题 17.4

1. (1) \({z}_{xx} = {12}{x}^{2} - 8{y}^{2},{z}_{xy} =  - {16xy},{z}_{yy} = {12}{y}^{2} - 8{x}^{2}\) ;


(2) \({z}_{xx} = {\mathrm{e}}^{x}\left( {\cos y + x\sin y + 2\sin y}\right) ,{z}_{xy} = {\mathrm{e}}^{x}\left( {x\cos y + \cos y - \sin y}\right) ,{z}_{yy} =  - {\mathrm{e}}^{x}\left( {\cos y + x\sin y}\right)\) ;

(3) \(\frac{{\partial }^{3}z}{\partial {x}^{2}\partial y} = 0,\frac{{\partial }^{3}z}{\partial x\partial {y}^{2}} =  - \frac{1}{{y}^{2}}\) ; (4) \(\frac{{\partial }^{p + q + r}u}{\partial {x}^{p}\partial {y}^{q}\partial {z}^{r}} = \left( {x + p}\right) \left( {y + q}\right) \left( {z + r}\right) {\mathrm{e}}^{x + y + z}\) ;

(5) \({z}_{xx} = {y}^{4}{f}_{11} + {4x}{y}^{3}{f}_{12} + 4{x}^{2}{y}^{2}{f}_{22} + {2y}{f}_{2},{z}_{xy} = {2x}{y}^{3}{f}_{11} + 5{x}^{2}{y}^{2}{f}_{12} + 2{x}^{3}y{f}_{22} + {2y}{f}_{1} + {2x}{f}_{2}\) ,

\({z}_{yy} = 4{x}^{2}{y}^{2}{f}_{11} + 4{x}^{3}y{f}_{12} + {x}^{4}{f}_{22} + {2x}{f}_{1};\)

( 6 ) \({u}_{xx} = 2{f}^{\prime } + 4{x}^{2}{f}^{\prime \prime },{u}_{yy} = 2{f}^{\prime } + 4{y}^{2}{f}^{\prime \prime },{u}_{zz} = 2{f}^{\prime } + 4{z}^{2}{f}^{\prime \prime },{u}_{xy} = {4xy}{f}^{\prime \prime },{u}_{yz} = {4yz}{f}^{\prime \prime },{u}_{xz} = {4xz}{f}^{\prime \prime }\) ;

(7) \({z}_{x} = {f}_{1} + y{f}_{2} + \frac{1}{y}{f}_{3},{z}_{xx} = {f}_{11} + {2y}{f}_{12} + \frac{2}{y}{f}_{13} + {y}^{2}{f}_{22} + 2{f}_{23} + \frac{1}{{y}^{2}}{f}_{33}\) ,

\({z}_{xy} = {f}_{11} + \left( {x + y}\right) {f}_{12} + \frac{1}{y}\left( {1 - \frac{x}{y}}\right) {f}_{13} + {xy}{f}_{22} - \frac{x}{{y}^{3}}{f}_{33} + {f}_{2} - \frac{1}{{y}^{2}}{f}_{3}.\)

7. (1) \({x}^{2} + {y}^{2} + {R}_{2},{R}_{2} =  - \frac{2}{3}\left\lbrack  {{3\theta }{\left( {x}^{2} + {y}^{2}\right) }^{2}\sin \left( {{\theta }^{2}{x}^{2} + {\theta }^{2}{y}^{2}}\right)  + 2{\theta }^{3}{\left( {x}^{2} + {y}^{2}\right) }^{3}\cos \left( {{\theta }^{2}{x}^{2} + {\theta }^{2}{y}^{2}}\right) }\right\rbrack\) ;

( 2 ) \(f\left( {x,y}\right)  = f\left( {1 + h,1 + k}\right)  = 1 + h - k - {hk} + {k}^{2} + h{k}^{2} - {k}^{3} + \left\lbrack  {-\frac{h{k}^{3}}{{\left( 1 + \theta k\right) }^{4}} + \frac{1 + {\theta h}}{{\left( 1 + \theta k\right) }^{5}}{k}^{4}}\right\rbrack\) ;

(3) \(\mathop{\sum }\limits_{{p = 1}}^{n}{\left( -1\right) }^{p - 1}\frac{{\left( x + y\right) }^{p}}{p} + {\left( -1\right) }^{n}\frac{{\left( x + y\right) }^{n + 1}}{n + 1} \cdot  \frac{1}{{\left( 1 + \theta x + \theta y\right) }^{n + 1}}\) ;

(4) \(5 + 2{\left( x - 1\right) }^{2} - \left( {x - 1}\right) \left( {y + 2}\right)  - {\left( y + 2\right) }^{2}\) .

8. (1)(a, a)为极大点; (2)(1,0)为极小点; (3) \(\left( {\frac{1}{2}, - 1}\right)\) 为极小点.

9. (1) 最大值 \(f\left( {2,0}\right)  = f\left( {-2,0}\right)  = 4\) ,最小值 \(f\left( {0,2}\right)  = f\left( {0, - 2}\right)  =  - 4\) ; (2) 最大值 \(f\left( {1,0}\right)  = f( - 1\) , \(0) = f\left( {0, - 1}\right)  = f\left( {0,1}\right)  = 1\) ,最小值 \(f\left( {0,0}\right)  = 0;\left( 3\right)\) 最大值 \(\frac{3}{2}\sqrt{3}\) ,最小值 0 .

10. 等边三角形.

11. \(\left( {\frac{8}{5},\frac{16}{5}}\right)\) .

12. \(\left( {\frac{\mathop{\sum }\limits_{{i = 1}}^{n}{x}_{i}}{n},\frac{\mathop{\sum }\limits_{{i = 1}}^{n}{y}_{i}}{n}}\right)\) .

19. (1) 0; (2) 3(z-y)(x-y)(x-z) 或 \(3\left\lbrack  {{z}^{2}\left( {y - x}\right)  + {x}^{2}\left( {z - y}\right)  + {y}^{2}\left( {x - z}\right) }\right\rbrack\) ; (3) 0.

20. \(f\left( {x,y,z}\right)  + \left( {{2Ax} + {Dy} + {Fz}}\right) h + \left( {{2By} + {Dx} + {Ez}}\right) k + \left( {{2Cz} + {Ey} + {Fx}}\right) l + f\left( {h,k,l}\right)\) .

\subsection{第十七章总练习题}

2. \({f}_{x}\left( {0,0}\right)  = 1,{f}_{y}\left( {0,0}\right)  =  - 1\) ,不可微.

5. \({\varphi }_{xx} = {6x} + 2\left( {a + e + k}\right)\) .

6. \(\frac{{\partial }^{3}\Phi }{\partial x\partial y\partial z} = \left| \begin{array}{lll} {f}_{1}^{\prime } & {f}_{2}^{\prime } & {f}_{3}^{\prime } \\  {g}_{1}^{\prime } & {g}_{2}^{\prime } & {g}_{3}^{\prime } \\  {h}_{1}^{\prime } & {h}_{2}^{\prime } & {h}_{3}^{\prime } \end{array}\right|\) .



\section{第十八章 隐函数定理及其应用}

\subsection{习题 18.1}

3. (1) \(\frac{\mathrm{d}y}{\mathrm{\;d}x} =  - \frac{{2y} + {12}{x}^{2}{y}^{3}}{x + 9{x}^{3}{y}^{2}}\) ; (2) \(\frac{\mathrm{d}y}{\mathrm{\;d}x} = \frac{x + y}{x - y}\left( {x \neq  y}\right)\) ; (3) \({z}_{x} = \frac{y{\mathrm{e}}^{-{xy}}}{2 - {\mathrm{e}}^{z}};{z}_{y} = \frac{x{\mathrm{e}}^{-{xy}}}{2 - {\mathrm{e}}^{z}}\) ;

(4) \({y}^{\prime } =  - \frac{y}{\sqrt{{a}^{2} - {y}^{2}}},{y}^{\prime \prime } = \frac{{a}^{2}y}{{\left( {a}^{2} - {y}^{2}\right) }^{2}};\left( 5\right) {z}_{x} = \frac{1 - x}{z - 2},{z}_{y} =  - \frac{1 + y}{z - 2}\) ;

(6) \({z}_{x} = \frac{{f}_{1} + {yz}{f}_{2}}{1 - {f}_{1} - {xy}{f}_{2}},{x}_{y} = \frac{{f}_{1} + {xz}{f}_{2}}{{f}_{1} + {yz}{f}_{2}},{y}_{z} = \frac{1 - \left( {{f}_{1} + {xy}{f}_{2}}\right) }{{f}_{1} + {xz}{f}_{2}}\) .

4. \(\frac{\mathrm{d}z}{\mathrm{\;d}x} = \frac{2\left( {{x}^{2} - {y}^{2}}\right) }{x - {2y}},\frac{{\mathrm{d}}^{2}z}{\mathrm{\;d}{x}^{2}} = \frac{{4x} - {2y}}{x - {2y}} + \frac{6x}{{\left( x - 2y\right) }^{3}}\) .

5. \({u}_{x} = 2\left( {x + \frac{z{x}^{2} - y{z}^{2}}{{xy} - {z}^{2}}}\right) ,{z}_{xx} = \frac{{2xz}\left( {{y}^{3} - {3xyz} + {x}^{3} + {z}^{3}}\right) }{{\left( xy - {z}^{2}\right) }^{3}} = 0,{u}_{xx} = 2\left\lbrack  {1 + {\left( \frac{{x}^{2} - {yz}}{{xy} - {z}^{2}}\right) }^{2}}\right\rbrack\) .

7. (1) \({z}_{x} = {z}_{y} =  - 1,{z}_{xx} = {z}_{xy} = {z}_{yy} = 0\) ;

(2) \({z}_{x} =  - \left( {1 + \frac{{F}_{1} + {F}_{2}}{{F}_{3}}}\right) ,{z}_{y} =  - \left( {1 + \frac{{F}_{2}}{{F}_{3}}}\right)\) ,

\({z}_{xx} =  - {F}_{3}^{-3}\left\lbrack  {{F}_{3}^{2}\left( {{F}_{11} + 2{F}_{12} + {F}_{22}}\right)  - 2{F}_{3}\left( {{F}_{1} + {F}_{2}}\right) \left( {{F}_{13} + {F}_{23}}\right)  + {\left( {F}_{1} + {F}_{2}\right) }^{2}{F}_{33}}\right\rbrack  .\)

9. \(f\) 在 1 的某邻域上具有连续的导函数,且 \({f}^{\prime }\left( 1\right)  \neq  0\) .

\subsection{习题 18.2}

2. (1) \(\frac{\mathrm{d}y}{\mathrm{\;d}x} =  - \frac{{2x} - a}{2y},\frac{\mathrm{d}z}{\mathrm{\;d}x} =  - \frac{a}{2z}\) ;

(2) \({u}_{x} = \frac{{2v} + {uy}}{{4uv} - {xy}},{v}_{x} = \frac{-x - 2{u}^{2}}{{4uv} - {xy}},{u}_{y} = \frac{-y - 2{v}^{2}}{{4uv} - {xy}},{v}_{y} = \frac{{2u} + {xv}}{{4uv} - {xy}}\) ;

(3) \({u}_{x} = \frac{-u\left( {{2vy}{g}_{2} - 1}\right) {f}_{1} - {f}_{2}{g}_{1}}{\left( {1 - x{f}_{1}}\right) \left( {1 - {2vy}{g}_{2}}\right)  - {f}_{2}{g}_{1}},{v}_{x} = \frac{u{f}_{1}{g}_{1} + x{f}_{1}{g}_{1} - {g}_{1}}{\left( {1 - x{f}_{1}}\right) \left( {1 - {2vy}{g}_{2}}\right)  - {f}_{2}{g}_{1}}\) .

3. (1) \(\frac{\partial u}{\partial x} = \frac{\sin v}{{\mathrm{e}}^{u}\left( {\sin v - \cos v}\right)  + 1},\frac{\partial v}{\partial x} = \frac{\cos v - {\mathrm{e}}^{u}}{u{\mathrm{e}}^{u}\left( {\sin v - \cos v}\right)  + u},\frac{\partial u}{\partial y} = \frac{-\cos v}{{\mathrm{e}}^{u}\left( {\sin v - \cos v}\right)  + 1}\) ,

\(\frac{\partial v}{\partial y} = \frac{{\mathrm{e}}^{u} + \sin v}{u{\mathrm{e}}^{u}\left( {\sin v - \cos v}\right)  + u};\left( 2\right) {z}_{x} =  - {3uv}.\)

4. \({\left. \mathrm{d}z\right| }_{\left( 1,1\right) } = 0\) .

5. (1) \(\frac{\partial z}{\partial u} = \frac{\partial z}{\partial v};\left( 2\right) \frac{{\partial }^{2}z}{\partial u\partial v} = \frac{1}{2u}\frac{\partial z}{\partial v}\) .

6. \(\frac{\partial u}{\partial x} = \frac{\partial f}{\partial x},\frac{\partial u}{\partial y} = \frac{\partial f}{\partial y} + \left( {\frac{\partial \left( {h,f}\right) }{\partial \left( {z,t}\right) }/\frac{\partial \left( {g,h}\right) }{\partial \left( {z,t}\right) }}\right) \frac{\partial g}{\partial y}\) .

10. (1) \(x = \frac{u}{{u}^{2} + {v}^{2} + {w}^{2}},y = \frac{v}{{u}^{2} + {v}^{2} + {w}^{2}},z = \frac{w}{{u}^{2} + {v}^{2} + {w}^{2}};\left( 2\right)  - \frac{1}{{r}^{6}}\) .

习题 18.3

1. \(\frac{y}{{y}_{0}^{1/3}} + \frac{x}{{x}_{0}^{1/3}} = {a}^{2/3}\) . 2. (1) \(\frac{x}{a} + \frac{z}{c} = 1,y = \frac{b}{2},{ax} - {cz} = \frac{1}{2}\left( {{a}^{2} - {c}^{2}}\right)\) ; (2) \(\frac{x - 1}{8} = \frac{y + 1}{10} = \frac{z - 2}{7},8\left( {x - 1}\right)  + {10}\left( {y + 1}\right)  + 7\left( {z - 2}\right)  = 0\) . 3. (1) \(- 2\left( {x - 1}\right)  + \left( {y - 1}\right)  + \left( {z - 2}\right)  = 0,\frac{x - 1}{-2} = y - 1 = z - 2\) ; (2) \(\frac{x}{a} + \frac{y}{b} + \frac{z}{c} = \sqrt{3},a\left( {x - \frac{a}{\sqrt{3}}}\right)  = b\left( {y - \frac{b}{\sqrt{3}}}\right)  = c\left( {z - \frac{c}{\sqrt{3}}}\right)\) . 5. \(x + {4y} + {6z} =  \pm  {21}\) . 6. \(\left( {-1,1, - 1}\right) ,\left( {-\frac{1}{3},\frac{1}{9}, - \frac{1}{27}}\right)\) . 7. \(\pm  \frac{16}{243}\) . 9. \(\lambda  =  \pm  \frac{abc}{3\sqrt{3}}\) . 10. \(x + y = \frac{1}{2}\left( {1 \pm  \sqrt{2}}\right)\) . 习题 18.4

1. (1) 极小值 \(f\left( {\frac{1}{2},\frac{1}{2}}\right)  = \frac{1}{2}\) ; (2) 极小值 \(f\left( {c,c,c,c}\right)  = {4c}\) ; (3) 极小值 \(f\left( {\frac{1}{\sqrt{6}},\frac{1}{\sqrt{6}}, - \frac{2}{\sqrt{6}}}\right)  =\)  \(f\left( {\frac{1}{\sqrt{6}}, - \frac{2}{\sqrt{6}},\frac{1}{\sqrt{6}}}\right)  = f\left( {-\frac{2}{\sqrt{6}},\frac{1}{\sqrt{6}},\frac{1}{\sqrt{6}}}\right)  =  - \frac{1}{3\sqrt{6}}\) ,极大值 \(f\left( {-\frac{1}{\sqrt{6}}, - \frac{1}{\sqrt{6}},\frac{2}{\sqrt{6}}}\right)  = f\left( {-\frac{1}{\sqrt{6}},\frac{2}{\sqrt{6}}, - \frac{1}{\sqrt{6}}}\right)  =\)  \(f\left( {\frac{2}{\sqrt{6}}, - \frac{1}{\sqrt{6}}, - \frac{1}{\sqrt{6}}}\right)  = \frac{1}{3\sqrt{6}}.\)

2. (1) 立方体; (2) 立方体.

3. \(\frac{\left| A{x}_{0} + B{y}_{0} + C{z}_{0} + D\right| }{\sqrt{{A}^{2} + {B}^{2} + {C}^{2}}}\) .

5. 当 \({x}_{i} = \frac{{a}_{i}}{{\left( \sum {a}_{k}^{2}\right) }^{1/2}},i = 1,2,\cdots ,n\) 时, \(f\) 取得最大值 \(\sqrt{\sum {a}_{k}^{2}}\) .

6. 当 \({x}_{i} = \frac{{a}_{i}}{\sum {a}_{k}^{2}},i = 1,2,\cdots ,n\) 时, \(f\) 取得最小值 \(\frac{1}{\sum {a}_{k}^{2}}\) .

\subsection{第十八章总练习题}

3. \(\frac{\mathrm{d}y}{\mathrm{\;d}x} =  - \frac{{f}_{x} + {f}_{z}{g}_{x}}{{f}_{y} + {f}_{z}{g}_{y}},\frac{\mathrm{d}z}{\mathrm{\;d}x} = \frac{{g}_{x}{f}_{y} - {g}_{y}{f}_{x}}{{f}_{y} + {f}_{z}{g}_{y}}\) .

5. \(\frac{\partial u}{\partial x} = \frac{\partial \left( {g,h}\right) }{\partial \left( {v,w}\right) }/J,\frac{\partial u}{\partial y} = \frac{\partial \left( {h,f}\right) }{\partial \left( {v,w}\right) }/J,\frac{\partial u}{\partial z} = \frac{\partial \left( {f,g}\right) }{\partial \left( {v,w}\right) }/J\) ,其中 \(J = \frac{\partial \left( {f,g,h}\right) }{\partial \left( {u,v,w}\right) }\) .

6. (1) \(\frac{\partial u}{\partial x} = \frac{{f}_{x} + {g}_{x} - {2x}}{{2u} - {f}_{u} - {g}_{u}},\frac{\partial u}{\partial y} = \frac{{g}_{y}}{{2u} - {g}_{u} - {f}_{u}}\) ; (2) \(\frac{\partial u}{\partial x} = \frac{{f}_{1}}{1 - {f}_{1} - y{f}_{2}},\frac{\partial u}{\partial y} = \frac{u{f}_{2}}{1 - {f}_{1} - y{f}_{2}}\) .

9. (1) 极大值 1,极小值-1; (2) 极大值 \(\sqrt{\frac{1}{8}}a\) ,极小值 \(- \sqrt{\frac{1}{8}}a\) .

10. \(\frac{\mathrm{d}y}{\mathrm{\;d}x} = \left( {{\psi }_{u} + {\psi }_{v}\frac{\mathrm{d}v}{\mathrm{\;d}u}}\right) /\left( {{\varphi }_{u} + {\varphi }_{v}\frac{\mathrm{d}v}{\mathrm{\;d}u}}\right)\) ;

\(\frac{{\mathrm{d}}^{2}y}{\mathrm{\;d}{x}^{2}} = \left\{  {\left\lbrack  {{\psi }_{uu} + {\psi }_{uv}\frac{\mathrm{d}v}{\mathrm{\;d}u} + \left( {{\psi }_{vu} + {\psi }_{vv}\frac{\mathrm{d}v}{\mathrm{\;d}u}}\right) \frac{\mathrm{d}v}{\mathrm{\;d}u} + {\psi }_{v}\frac{{\mathrm{d}}^{2}v}{\mathrm{\;d}{u}^{2}}}\right\rbrack  \left( {{\varphi }_{u} + {\varphi }_{v}\frac{\mathrm{d}v}{\mathrm{\;d}u}}\right)  - }\right.\)

\(\left. {\left( {{\psi }_{u} + {\psi }_{v}\frac{\mathrm{d}v}{\mathrm{\;d}u}}\right) \left\lbrack  {{\varphi }_{uu} + {\varphi }_{uv}\frac{\mathrm{d}v}{\mathrm{\;d}u} + \left( {{\varphi }_{vu} + {\varphi }_{vv}\frac{\mathrm{d}v}{\mathrm{\;d}u}}\right) \frac{\mathrm{d}v}{\mathrm{\;d}u} + {\varphi }_{v}\frac{{\mathrm{d}}^{2}v}{\mathrm{\;d}{u}^{2}}}\right\rbrack  }\right\}  /{\left( {\varphi }_{u} + {\varphi }_{v}\frac{\mathrm{d}v}{\mathrm{\;d}u}\right) }^{3}.\)

13. \(\frac{\sqrt{3}}{2}{abc}\) .

\section{第十九章 含参量积分}

习题 19.1

1. \(F\left( y\right)  = \left\{  \begin{array}{ll} 1, &  - \infty  < y < 0, \\  1 - {2y}, & 0 \leq  y \leq  1, \\   - 1, & 1 < y <  + \infty . \end{array}\right.\)

2. (1) 1 ; (2) \(\frac{8}{3}\) .

3. \({\int }_{x}^{{x}^{2}} - {y}^{2}{\mathrm{e}}^{-x{y}^{2}}\mathrm{\;d}y + {2x}{\mathrm{e}}^{-{x}^{5}} - {\mathrm{e}}^{-{x}^{3}}\) .

4. (1) \(\pi \ln \frac{\left| a\right|  + \left| b\right| }{2}\) ; (2) \(\left\{  \begin{array}{ll} 0, & \left| a\right|  \leq  1, \\  {2\pi }\ln \left| a\right| , & \left| a\right|  > 1. \end{array}\right.\)

5. (1) \(\arctan \left( {1 + b}\right)  - \arctan \left( {1 + a}\right) ;\;\left( 2\right) \frac{1}{2}\ln \left( \frac{{b}^{2} + {2b} + 2}{{a}^{2} + {2a} + 2}\right)\) .

6. \(\frac{\pi }{4}, - \frac{\pi }{4}\) .

9. \(x\left( {2 - 3{y}^{2}}\right) f\left( {xy}\right)  + \frac{x}{{y}^{2}}f\left( \frac{x}{y}\right)  + {x}^{2}y\left( {1 - {y}^{2}}\right) {f}^{\prime }\left( {xy}\right)\) .

10. (1) \({E}^{\prime }\left( k\right)  = \frac{1}{k}\left\lbrack  {E\left( k\right)  - F\left( k\right) }\right\rbrack  ,{F}^{\prime }\left( k\right)  = \frac{E\left( k\right) }{k\left( {1 - {k}^{2}}\right) } - \frac{F\left( k\right) }{k}\) .

习题 19.2

2. \(\ln \frac{b}{a}\) .

4. (1) \(\sqrt{\pi }\left( {\left| b\right|  - \left| a\right| }\right)\) ; (3) yarctan \(y - \frac{1}{2}\ln \left( {1 + {y}^{2}}\right)\) .

7. \(\frac{\pi }{2}\frac{\left( {{2n} - 1}\right) !!}{\left( {2n}\right) !!}{a}^{-{2n} - 1}\) .

\subsection{习题 19.3}

1. \(\frac{3}{4}\sqrt{\pi }, - \frac{8}{15}\sqrt{\pi },\frac{\left( {{2n} - 1}\right) !!}{{2}^{n}}\sqrt{\pi },\frac{{\left( -1\right) }^{n}{2}^{n}}{\left( {{2n} - 1}\right) !!}\sqrt{\pi }\) .

2. \(\frac{\left( {{2n} - 1}\right) !!}{\left( {2n}\right) !!} \cdot  \frac{\pi }{2},\frac{\left( {2n}\right) !!}{\left( {{2n} + 1}\right) !!}\) .

6. (1) \(\frac{1}{2}\mathrm{\;B}\left( {\frac{m + 1}{2},\frac{n + 1}{2}}\right) ,m >  - 1,n >  - 1\) ;

(2) \(\Gamma \left( {p + 1}\right) ,p >  - 1\) .

\subsection{第十九章总练习题}

1. \(a =  - \frac{11}{3},b = 4\) .

3. \(F\left( a\right)  = \frac{\pi }{2}\operatorname{sgn}\left( {1 - {a}^{2}}\right) ,a =  \pm  1\) .

\section{第二十章 曲 线 积 分}

\subsection{习题 20.1}

1. (1) \(1 + \sqrt{2}\) ; (3) \(\frac{{ab}\left( {{a}^{2} + {ab} + {b}^{2}}\right) }{3\left( {a + b}\right) }\) ; (4) 4 ;

(5) \(\sqrt{{a}^{2} + {b}^{2}}\left( {2{a}^{2}\pi  + \frac{8{\pi }^{3}{b}^{2}}{3}}\right)\) ； (6) \(\frac{{16}\sqrt{2}}{143};\) (7) \({2\pi }{a}^{2}\) .

2. \(\frac{a}{3}\left( {2\sqrt{2} - 1}\right)\) .

3. \(\bar{x} = \frac{4}{3}a,\bar{y} = \frac{4}{3}a\) .

4. (1) \(\frac{\pi }{4}a{\mathrm{e}}^{a}\) ; (2) \(\frac{{2k}{a}^{2}\sqrt{{k}^{2} + 1}}{4{k}^{2} + 1}\) .

\subsection{习题 20.2}

1. (1) \(\frac{2}{3},0,2\) ; (2) \({a}^{2}\pi\) ; (3) 0; (4) 2; (5) 13.

2. \(\frac{k}{2}\left( {{b}^{2} - {a}^{2}}\right) ,k\) 为比例系数.

3. \(\frac{-k\sqrt{{a}^{2} + {b}^{2} + {c}^{2}}}{\left| c\right| }\ln 2,k\) 为比例系数.

5. (1) \(\frac{\sqrt{2}}{16}\pi\) ; (2) -4 .

\subsection{第二十章总练习题}

1. (1) \(\frac{1}{12}\left\lbrack  {5\sqrt{5} - {17}\sqrt{17}}\right\rbrack   - \frac{3}{2}\sqrt{2}\) ; (2) \(4{a}^{2}\left( {1 - \frac{\sqrt{2}}{2}}\right)\) ;

(3) \(\frac{1}{3}\left\lbrack  {{\left( 2 + {t}_{0}\right) }^{\frac{3}{2}} - 2\sqrt{2}}\right\rbrack\) ; (4) \(- \pi \frac{{a}^{4}}{4}\) ; (5) \(\ln 2\) ; (6) \(- \frac{\pi }{4}{a}^{3}\) .

2. (1) \({\int }_{a}^{b}f\left( {x,a}\right) \mathrm{d}x,{\int }_{a}^{b}f\left( {x,a}\right) \mathrm{d}x,0\) ;

(2) \({\int }_{a}^{b}f\left( {x,a}\right) \mathrm{d}x + {\int }_{a}^{b}f\left( {b,y}\right) \mathrm{d}y + {\int }_{a}^{b}\sqrt{2}f\left( {t,t}\right) \mathrm{d}t\) ,

\({\int }_{a}^{b}f\left( {x,a}\right) \mathrm{d}x + {\int }_{b}^{a}f\left( {t,t}\right) \mathrm{d}t,{\int }_{a}^{b}f\left( {b,y}\right) \mathrm{d}y + {\int }_{b}^{a}f\left( {t,t}\right) \mathrm{d}t.\)

\section{第二十一章 重 积 分}

习题 21.1

1. \(\frac{1}{4}\) .

8. \(\frac{200}{102} \leq  I \leq  2\) .

习题 21.2

1. (1) \({\int }_{a}^{b}\mathrm{\;d}y{\int }_{y}^{b}f\left( {x,y}\right) \mathrm{d}x = {\int }_{a}^{b}\mathrm{\;d}x{\int }_{a}^{x}f\left( {x,y}\right) \mathrm{d}y\) ;

(2) \({\int }_{0}^{\sqrt{2}}\mathrm{\;d}y{\int }_{y}^{\sqrt{1 - {y}^{2}}}f\left( {x,y}\right) \mathrm{d}x = {\int }_{0}^{\sqrt{2}}\mathrm{\;d}x{\int }_{0}^{x}f\left( {x,y}\right) \mathrm{d}y + {\int }_{\frac{\sqrt{2}}{2}}^{1}\mathrm{\;d}x{\int }_{0}^{\sqrt{1 - {x}^{2}}}f\left( {x,y}\right) \mathrm{d}y\) ;

(3) \({\int }_{0}^{1}\mathrm{\;d}x{\int }_{1 - x}^{\sqrt{1 - {x}^{2}}}f\left( {x,y}\right) \mathrm{d}y = {\int }_{0}^{1}\mathrm{\;d}y{\int }_{1 - y}^{\sqrt{1 - {y}^{2}}}f\left( {x,y}\right) \mathrm{d}x\) ;

(4) \({\int }_{-1}^{0}\mathrm{\;d}x{\int }_{-1 - x}^{x + 1}f\left( {x,y}\right) \mathrm{d}y + {\int }_{0}^{1}\mathrm{\;d}x{\int }_{x - 1}^{1 - x}f\left( {x,y}\right) \mathrm{d}y\)

\(= {\int }_{0}^{1}\mathrm{\;d}y{\int }_{y - 1}^{1 - y}f\left( {x,y}\right) \mathrm{d}x + {\int }_{-1}^{0}\mathrm{\;d}y{\int }_{-1 - y}^{y + 1}f\left( {x,y}\right) \mathrm{d}x.\)

2. (1) \({\int }_{0}^{2}\mathrm{\;d}y{\int }_{\frac{y}{2}}^{y}f\left( {x,y}\right) \mathrm{d}x + {\int }_{2}^{4}\mathrm{\;d}y{\int }_{\frac{y}{2}}^{2}f\left( {x,y}\right) \mathrm{d}x\) ;

( 2 ) \({\int }_{-1}^{0}\mathrm{\;d}y{\int }_{-\sqrt{1 - {y}^{2}}}^{\sqrt{1 - {y}^{2}}}f\left( {x,y}\right) \mathrm{d}x + {\int }_{0}^{1}\mathrm{\;d}y{\int }_{-\sqrt{1 - y}}^{\sqrt{1 - y}}f\left( {x,y}\right) \mathrm{d}x\) ；

(3) \({\int }_{0}^{a}\mathrm{\;d}y{\int }_{\frac{{y}^{2}}{2a}}^{a - \sqrt{{a}^{2} - {y}^{2}}}f\left( {x,y}\right) \mathrm{d}x + {\int }_{0}^{a}\mathrm{\;d}y{\int }_{a + \sqrt{{a}^{2} - {y}^{2}}}^{2a}f\left( {x,y}\right) \mathrm{d}x + {\int }_{a}^{2a}\mathrm{\;d}y{\int }_{\frac{{y}^{2}}{2a}}^{2a}f\left( {x,y}\right) \mathrm{d}x\) ;

(4) \({\int }_{0}^{1}\mathrm{\;d}y{\int }_{\sqrt{y}}^{3 - {2y}}f\left( {x,y}\right) \mathrm{d}x\) .

3. (1) \(\frac{{p}^{5}}{21}\) ; (2) \(\frac{128}{105}\) ; (3) \(\left( {2\sqrt{2} - \frac{8}{3}}\right) {a}^{\frac{3}{2}}\) ; (4) \(\frac{8}{15}\) .

4. \(9\frac{1}{6}\) .

习题 21.3

1. (1) \(- {46}\frac{2}{3}\) ; (2) \(\frac{m{a}^{2}\pi }{8}\) .

2. (1) \(\frac{3}{8}{a}^{2}\pi\) ; (2) \({a}^{2}\) .

4. \({2\sigma },\sigma\) 为由 \(L\) 所围的面积.

5. (1) 0; (2) \({y}^{2}\cos x + {x}^{2}\cos y;\left( 3\right)  - \frac{3}{2};\left( 4\right) 9;\left( 5\right) {\int }_{2}^{1}\varphi \left( x\right) \mathrm{d}x + {\int }_{1}^{2}\psi \left( y\right) \mathrm{d}y\) .

6. (1) \(\frac{1}{3}{x}^{3} + {x}^{2}y - x{y}^{2} - \frac{1}{3}{y}^{3} + c\) ; (2) \({\mathrm{e}}^{x + y}\left( {x - y + 1}\right)  + y{\mathrm{e}}^{x} + c\) ;

(3) \(\frac{1}{2}\int f\left( \sqrt{u}\right) \mathrm{d}u,u = {x}^{2} + {y}^{2}\) .

7. \(y{F}_{y}\left( {x,y}\right)  = x{F}_{x}\left( {x,y}\right)\) .

8. \(\pm  {mS} + \varphi \left( {y}_{2}\right) {\mathrm{e}}^{{x}_{2}} - \varphi \left( {y}_{1}\right) {\mathrm{e}}^{{x}_{1}} - \frac{m}{2}\left( {{x}_{2} - {x}_{1}}\right) \left( {{y}_{1} + {y}_{2}}\right)  - m\left( {{y}_{2} - {y}_{1}}\right)\) .

习题 21.4

1. (1) \({\int }_{0}^{\pi }\mathrm{d}\theta {\int }_{a}^{b}{rf}\left( {r\cos \theta ,r\sin \theta }\right) \mathrm{d}r = {\int }_{a}^{b}\mathrm{\;d}r{\int }_{0}^{\pi }{rf}\left( {r\cos \theta ,r\sin \theta }\right) \mathrm{d}\theta\) ;

(2) \({\int }_{0}^{\frac{\pi }{2}}\mathrm{\;d}\theta {\int }_{0}^{\sin \theta }{rf}\left( {r\cos \theta ,r\sin \theta }\right) \mathrm{d}r = {\int }_{0}^{1}\mathrm{\;d}r{\int }_{\arcsin r}^{\frac{\pi }{2}}{rf}\left( {r\cos \theta ,r\sin \theta }\right) \mathrm{d}\theta\) ;

(3) \({\int }_{-\frac{\pi }{4}}^{0}\mathrm{\;d}\theta {\int }_{0}^{\sec \theta }{rf}\left( {r\cos \theta ,r\sin \theta }\right) \mathrm{d}r + {\int }_{0}^{\frac{\pi }{2}}\mathrm{\;d}\theta {\int }_{0}^{\frac{1}{\cos \theta  + \sin \theta }}{rf}\left( {r\cos \theta ,r\sin \theta }\right) \mathrm{d}r\)

\(= {\int }_{0}^{\frac{\sqrt{2}}{2}}\mathrm{\;d}r{\int }_{-\frac{\pi }{4}}^{\frac{\pi }{2}}{rf}\left( {r\cos \theta ,r\sin \theta }\right) \mathrm{d}\theta  + {\int }_{\frac{\sqrt{2}}{2}}^{1}\mathrm{\;d}r{\int }_{-\frac{\pi }{4}}^{\frac{\pi }{4} - \arccos \frac{1}{\sqrt{2}}r}{rf}\left( {r\cos \theta ,r\sin \theta }\right) \mathrm{d}\theta  +\)

\({\int }_{\frac{\sqrt{2}}{2}}^{1}\mathrm{\;d}r{\int }_{\frac{\pi }{4} + \arccos \frac{1}{\sqrt{2}r}}^{\frac{\pi }{2}}{rf}\left( {r\cos \theta ,r\sin \theta }\right) \mathrm{d}\theta  + {\int }_{1}^{\sqrt{2}}\mathrm{\;d}r{\int }_{-\frac{\pi }{4}}^{-\arccos \frac{1}{r}}{rf}\left( {r\cos \theta ,r\sin \theta }\right) \mathrm{d}\theta .\)

2. (1) \(- 6{\pi }^{2};\;\left( 2\right) \frac{\pi }{2}\) ; (3) \(\frac{{a}^{4}}{2}\) ; (4) \(\pi \left\lbrack  {f\left( {R}^{2}\right)  - f\left( 0\right) }\right\rbrack\) .

3. (1) \({\int }_{1}^{2}\mathrm{\;d}u{\int }_{-u}^{4 - u}\frac{1}{2}f\left( {\frac{u + v}{2},\frac{u - v}{2}}\right) \mathrm{d}v\) ;

(2) \({\int }_{0}^{\frac{\pi }{2}}\mathrm{\;d}v{\int }_{0}^{a}f\left( {u{\cos }^{4}v,u{\sin }^{4}v}\right) {4u}{\sin }^{3}v{\cos }^{3}v\mathrm{\;d}u\) ;

(3) \({\int }_{0}^{a}\mathrm{\;d}u{\int }_{0}^{1}f\left( {u\left( {1 - v}\right) ,{uv}}\right) u\mathrm{\;d}v\) .

4. (1) \(u = x + y,v = x - y,\frac{{\pi }^{2}}{2};\left( 2\right) u = x + y,v = y,\frac{\mathrm{e} - 1}{2}\) .

5. (1) \(\frac{\pi }{8}\) ; (2) \({8\pi }\) .

6. (1) \(\frac{{b}^{2} - {a}^{2}}{2}\left( {\frac{1}{1 + \alpha } - \frac{1}{1 + \beta }}\right)\) ; (2) \(\frac{{ab}\left( {{a}^{2} + {b}^{2}}\right) \pi }{2}\) ; (3) \(\left( {\sqrt{3} - \frac{\pi }{3}}\right) {a}^{2}\) .

8. (1) 极坐标变换, \({2\pi }{\int }_{0}^{1}{rf}\left( r\right) \mathrm{d}r\) ;

(2)极坐标变换, \(\pi {\int }_{0}^{\sqrt{2}}{rf}\left( r\right) \mathrm{d}r - 4{\int }_{1}^{\sqrt{2}}r\arccos \frac{1}{r}f\left( r\right) \mathrm{d}r\) ;

(3) \(u = x + y,v = x - y,{\int }_{-1}^{1}f\left( u\right) \mathrm{d}u\) ;

(4) \(u = {xy},v = \frac{y}{x},\ln 2{\int }_{1}^{2}f\left( u\right) \mathrm{d}u\) .

\subsection{习题 21.5}

1. (1) 14; (2) \(\frac{1}{2};\;\left( 3\right) \frac{1}{2}\left( {\ln 2 - \frac{5}{8}}\right) ;\;\left( 4\right) \frac{{\pi }^{2}}{16} - \frac{1}{2}\) .

2. (1) \(I = {\int }_{0}^{1}\mathrm{\;d}y{\int }_{0}^{1 - y}\mathrm{\;d}x{\int }_{0}^{x + y}f\left( {x,y,z}\right) \mathrm{d}z\)

\[
= {\int }_{0}^{1}\mathrm{\;d}x{\int }_{0}^{x}\mathrm{\;d}z{\int }_{0}^{1 - x}f\left( {x,y,z}\right) \mathrm{d}y + {\int }_{0}^{1}\mathrm{\;d}x{\int }_{x}^{1}\mathrm{\;d}z{\int }_{z - x}^{1 - x}f\left( {x,y,z}\right) \mathrm{d}y
\]

\(= {\int }_{0}^{1}\mathrm{\;d}z{\int }_{z}^{1}\mathrm{\;d}x{\int }_{0}^{1 - x}f\left( {x,y,z}\right) \mathrm{d}y + {\int }_{0}^{1}\mathrm{\;d}z{\int }_{0}^{z}\mathrm{\;d}x{\int }_{z - x}^{1 - x}f\left( {x,y,z}\right) \mathrm{d}y\)

\[
= {\int }_{0}^{1}\mathrm{\;d}y{\int }_{0}^{y}\mathrm{\;d}z{\int }_{0}^{1 - y}f\left( {x,y,z}\right) \mathrm{d}x + {\int }_{0}^{1}\mathrm{\;d}y{\int }_{y}^{1}\mathrm{\;d}z{\int }_{z - y}^{1 - y}f\left( {x,y,z}\right) \mathrm{d}x
\]

\[
= {\int }_{0}^{1}\mathrm{\;d}z{\int }_{z}^{1}\mathrm{\;d}y{\int }_{0}^{1 - y}f\left( {x,y,z}\right) \mathrm{d}x + {\int }_{0}^{1}\mathrm{\;d}z{\int }_{0}^{z}\mathrm{\;d}y{\int }_{z - y}^{1 - y}f\left( {x,y,z}\right) \mathrm{d}x;
\]

(2) \(I = {\int }_{0}^{1}\mathrm{\;d}y{\int }_{0}^{1}\mathrm{\;d}x{\int }_{0}^{{x}^{2} + {y}^{2}}f\left( {x,y,z}\right) \mathrm{d}z\)

\[
= {\int }_{0}^{1}\mathrm{\;d}x{\int }_{0}^{{x}^{2}}\mathrm{\;d}z{\int }_{0}^{1}f\left( {x,y,z}\right) \mathrm{d}y + {\int }_{0}^{1}\mathrm{\;d}x{\int }_{{x}^{2}}^{{x}^{2} + 1}\mathrm{\;d}z{\int }_{\sqrt{z - {x}^{2}}}^{1}f\left( {x,y,z}\right) \mathrm{d}y
\]

\[
= {\int }_{0}^{1}\mathrm{\;d}z{\int }_{\sqrt{z}}^{1}\mathrm{\;d}x{\int }_{0}^{1}f\left( {x,y,z}\right) \mathrm{d}y + {\int }_{0}^{1}\mathrm{\;d}z{\int }_{0}^{\sqrt{z}}\mathrm{\;d}x{\int }_{\sqrt{z - {x}^{2}}}^{1}f\left( {x,y,z}\right) \mathrm{d}y +
\]

\[
{\int }_{1}^{2}\mathrm{\;d}z{\int }_{\sqrt{z - 1}}^{1}\mathrm{\;d}x{\int }_{\sqrt{z - {x}^{2}}}^{1}f\left( {x,y,z}\right) \mathrm{d}y
\]

\[
= {\int }_{0}^{1}\mathrm{\;d}y{\int }_{0}^{{y}^{2}}\mathrm{\;d}z{\int }_{0}^{1}f\left( {x,y,z}\right) \mathrm{d}x + {\int }_{0}^{1}\mathrm{\;d}y{\int }_{{y}^{2}}^{{y}^{2} + 1}\mathrm{\;d}z{\int }_{\sqrt{z - {y}^{2}}}^{1}f\left( {x,y,z}\right) \mathrm{d}x
\]

\[
= {\int }_{0}^{1}\mathrm{\;d}z{\int }_{\sqrt{z}}^{1}\mathrm{\;d}y{\int }_{0}^{1}f\left( {x,y,z}\right) \mathrm{d}x + {\int }_{0}^{1}\mathrm{\;d}z{\int }_{0}^{\sqrt{z}}\mathrm{\;d}y{\int }_{\sqrt{z - {y}^{2}}}^{1}f\left( {x,y,z}\right) \mathrm{d}x +
\]

\[
{\int }_{1}^{2}\mathrm{\;d}z{\int }_{\sqrt{z - 1}}^{1}\mathrm{\;d}y{\int }_{\sqrt{z - {y}^{2}}}^{1}f\left( {x,y,z}\right) \mathrm{d}x.
\]

3. (1) \(\frac{59}{480}\pi {r}^{5};\;\left( 2\right) \frac{\pi }{15}\left( {2\sqrt{2} - 1}\right)\) .

4. (1) 柱坐标变换, \(\frac{3}{35}\) ;

(2) \(x = {ar}\sin \varphi {\cos }^{2}\theta ,y = {br}\sin \varphi {\sin }^{2}\theta ,z = {cr}\cos \varphi ,\frac{1}{3}{abc}\) .

5. \(\frac{8}{5}\pi\) .

7. (1) \(\frac{1}{4}{abc}{\pi }^{2}\) ; (2) \({4\pi abc}\left( {\mathrm{e} - 2}\right)\) .

习题 21.6

1. \(\frac{{2\pi }{a}^{2}}{3}\left( {2\sqrt{2} - 1}\right)\) .

2. \(\sqrt{2}\pi\) .

3. (1) \(\bar{x} = 0,\bar{y} = \frac{4b}{3\pi }\) ; (2) \(\bar{x} = 0,\bar{y} = \frac{\left( 2b + a\right) }{3\left( {a + b}\right) }h\) .

4. (1) \(\bar{x} = \bar{y} = 0,\bar{z} = \frac{1}{3}\) ; (2) \(\bar{x} = \frac{1}{4},\bar{y} = \frac{1}{8},\bar{z} =  - \frac{1}{4}\) .

5. (1) \(\frac{5\pi }{4}\rho {R}^{4}\) ; (2) \(\frac{1}{3}\rho {a}^{3}b{\sin }^{3}\varphi\) .

6. (1) \(\left( {0,0, - {2k\pi \rho }\left( {1 - \frac{c}{\sqrt{{R}^{2} + {c}^{2}}}}\right) }\right)\) ;

( 2 ) \(\left( {0,0,{2k\pi \rho }\left( {\sqrt{{a}^{2} + {c}^{2}} - \sqrt{{a}^{2} + \left( {h - {c}^{2}}\right) } - h}\right) }\right)\) ;

(3) \(\left( {0,0,\frac{{2k\pi \rho mh}\left( {h - \sqrt{{R}^{2} + {h}^{2}}}\right) }{\sqrt{{R}^{2} + {h}^{2}}}}\right)\) .

7. \({4ab}{\pi }^{2}\) .

8. \(\pi \left\lbrack  {a\sqrt{{a}^{2} + {b}^{2}} + {b}^{2}\ln \left( {a + \sqrt{{a}^{2} + {b}^{2}}}\right)  - {b}^{2}\ln b}\right\rbrack\) .

9. \(\frac{2}{3}\rho {a}^{5}\) .

习题 21.7

1. \(\frac{8}{15}{\pi }^{2}{r}^{5}\) .

2. \({\pi }^{2}\left( {1 - \frac{\pi }{4}}\right)\) .

3. \(\frac{1}{n!}{a}_{1}{a}_{2}\cdots {a}_{n}\) .

4. \(\frac{2{\pi }^{\frac{n}{2}}}{\Gamma \left( \frac{n}{2}\right) }{\int }_{0}^{R}{r}^{n - 1}f\left( r\right) \mathrm{d}r\) .

\subsection{习题 21.8}

1. (1) \(m > 1\) 收敛; (2) \(p > 1,q > 1\) 收敛; (3) \(p > \frac{1}{2}\) 收敛.

2. \(\frac{\pi }{2}\) .

3. (1) \(m < 1\) 收敛; (2) \(m < 1\) 收敛.

\subsection{第二十一章总练习题}

1. (1) \(\frac{1}{4}\) ; (2) \(\frac{6}{5}\) .

2. (1) 6; (2) 4 [ \(\frac{\pi }{3} + \ln \left( {2 + \sqrt{3}}\right)\) ].

3. \(\frac{\pi }{4}{a}^{4}\) .

4. \(f\left( {0,0}\right)\) .

5. (1) \(\frac{2}{t}F\left( t\right) ;\;\left( 2\right) {4\pi }{t}^{2}f\left( {t}^{2}\right)\) ;

(3) \(\frac{3}{t}\left\lbrack  {F\left( t\right)  + {\iint }_{\begin{matrix} {0 \leq  x \leq  t} \\  {0 \leq  y \leq  t} \\  {0 \leq  z \leq  t} \end{matrix}}{xyz}{f}^{\prime }\left( {xyz}\right) \mathrm{d}x\mathrm{\;d}y\mathrm{\;d}z}\right\rbrack\) ,其中 \(t > 0\) .

6. \(\frac{1}{4}\left( {{\mathrm{e}}^{-1} - 1}\right)\) .

8. 柱面坐标系: \({\int }_{0}^{1}\mathrm{\;d}z{\int }_{0}^{\frac{\pi }{4}}\mathrm{\;d}\theta {\int }_{0}^{\frac{1}{\cos \theta }}{rf}\left( {r\cos \theta ,r\sin \theta ,z}\right) \mathrm{d}r + {\int }_{0}^{1}\mathrm{\;d}z{\int }_{\frac{\pi }{4}}^{\frac{\pi }{2}}\mathrm{\;d}\theta {\int }_{0}^{\frac{1}{\sin \theta }}{rf}\left( {r\cos \theta ,r\sin \theta ,z}\right) \mathrm{d}r\) ; 球面坐标系: \({\int }_{0}^{\frac{\pi }{4}}\mathrm{\;d}\theta {\int }_{0}^{\operatorname{arccot}\cos \theta }\mathrm{d}\varphi {\int }_{0}^{\frac{1}{\cos \varphi }}{kf}\left( {u,v,w}\right) \mathrm{d}r + {\int }_{0}^{\frac{\pi }{4}}\mathrm{\;d}\theta {\int }_{\operatorname{arccot}\cos \theta }^{\frac{1}{2}}\mathrm{\;d}\varphi {\int }_{0}^{\frac{1}{\sin \varphi \cos \theta }}{kf}\left( {u,v,w}\right) \mathrm{d}r +\)

\({\int }_{\frac{\pi }{4}}^{\frac{\pi }{2}}\mathrm{\;d}\theta {\int }_{0}^{\operatorname{arccot}\sin \theta }\mathrm{d}\varphi {\int }_{0}^{\frac{1}{\cos \varphi }}{kf}\left( {u,v,w}\right) \mathrm{d}r + {\int }_{\frac{\pi }{4}}^{\frac{\pi }{2}}\mathrm{\;d}\theta {\int }_{\operatorname{arccot}\sin \theta }^{\frac{\pi }{2}}\mathrm{\;d}\varphi {\int }_{0}^{\frac{1}{\sin \varphi \sin \theta }}{kf}\left( {u,v,w}\right) \mathrm{d}r,\)

其中 \(k = {r}^{2}\sin \varphi ,u = r\sin \varphi \cos \theta ,v = r\sin \varphi \sin \theta ,w = r\cos \varphi\) .

10. \({2\pi }\) .

11. \(\frac{\pi }{\left| {a}_{1}{b}_{2} - {a}_{2}{b}_{1}\right| }\) .

12. \(V = \frac{8}{\left| \Delta \right| }{h}_{1}{h}_{2}{h}_{3},\Delta  = \left| \begin{array}{lll} {a}_{1} & {b}_{1} & {c}_{1} \\  {a}_{2} & {b}_{2} & {c}_{2} \\  {a}_{3} & {b}_{3} & {c}_{3} \end{array}\right|\) .

13. \(\left( {-\frac{2amk}{r},\frac{ma\pi k}{r}}\right) ,k\) 为引力常数.

14. \({2\pi }{a}^{2}\sqrt{{a}^{2} + {b}^{2}}\) .

15. \(\bar{x} = \frac{4}{3}a,\bar{y} = \frac{4}{3}a\) .

17. \(\lambda  =  - 1,\frac{\sqrt{{s}^{2} + {t}^{2}}}{t} - \frac{\sqrt{{s}_{0}^{2} + {t}_{0}^{2}}}{{t}_{0}}\) .

\section{第二十二章 曲面积分}

\subsection{习题 22.1}

(2) \(\frac{\pi }{2}\left( {\sqrt{2} + 1}\right)\) ; (3) \(\frac{2\pi H}{R}\) ; (4) \(\frac{\sqrt{3}}{120}\) .

2. \(\left( {\frac{a}{2},\frac{a}{2},\frac{a}{2}}\right)\) .

3. \(\frac{4}{3}{\pi \rho }{a}^{4}\) .

4. \(\frac{\pi {a}^{4}}{2}\sin \theta {\cos }^{2}\theta\) .

\subsection{习题 22.2}

(3) \(\frac{1}{8}\) ; (4) \(\frac{\pi }{4}\) ; (5) \(\frac{3}{8}\pi {R}^{3}\left( {a + b + c}\right)\) .

2. \(\frac{32}{3}\pi\) .

3. \(I = {bc}\left\lbrack  {f\left( a\right)  - f\left( 0\right) }\right\rbrack   + {ac}\left\lbrack  {g\left( b\right)  - g\left( 0\right) }\right\rbrack   + {ab}\left\lbrack  {h\left( c\right)  - h\left( 0\right) }\right\rbrack\) .

4. \({2\pi }{a}^{3}\) .

习题 22.3

1. (1) 0 ; (2) \(3{a}^{4}\) ; (3) \(\frac{\pi }{2}{h}^{4}\) ; (4) \(\frac{12}{5}\pi\) ; (5) \({2\pi }{a}^{3}\) .

2. \(\frac{11}{24}\) .

3. (1) 0; (2) 0; (3) 3 \({a}^{2}\) .

4. (1) \({xyz} + c;\;\left( 2\right) \frac{1}{3}\left( {{x}^{3} + {y}^{3} + {z}^{3}}\right)  - {2xyz} + c\) .

5. (1) \(- {53}\frac{7}{12}\) ; (2) 0 .

9. \({2S}\) .

\subsection{习题 22.4}

1. \(\frac{1}{r}\left( {x,y,z}\right) ,2\left( {x,y,z}\right) , - \frac{1}{{r}^{3}}\left( {x,y,z}\right) ,{f}^{\prime }\left( r\right) \frac{1}{r}\left( {x,y,z}\right) ,n{r}^{n - 2}\left( {x,y,z}\right)\) .

2. \(\left( {-4,2, - 4}\right) ,\left( {0,8,2}\right) ,\left( {-8, - 4, - {10}}\right) ,\nabla u\left( {5, - 3,\frac{2}{3}}\right)  = 0\) .

4. (1) \(0,2\left( {y - z,z - x,x - y}\right) ;\;\left( 2\right) {6xyz},\left( {x\left( {{z}^{2} - {y}^{2}}\right) ,y\left( {{x}^{2} - {z}^{2}}\right) ,z\left( {{y}^{2} - {x}^{2}}\right) }\right)\) ;

(3) \(\frac{x + y + z}{xyz},\frac{1}{xyz}\left( {\frac{{y}^{2}}{z} - \frac{{z}^{2}}{y},\frac{{z}^{2}}{x} - \frac{{x}^{2}}{z},\frac{{x}^{2}}{y} - \frac{{y}^{2}}{x}}\right)\) .

8. \(\frac{3}{8}\pi\) .

9. (1) \({2\pi }\) ; (2) \({2\pi }\) .

\subsection{第二十二章总练习题}

1. (1) \(\pi {a}^{2}\left( {1 - \lambda }\right) \left( {5 - {3\pi c}}\right)  - 8{\pi }^{2}{c}^{2}\) ; (2) rot \(\mathbf{A} = \left( {2\left( {1 - \lambda }\right) x,\left( {1 - \lambda }\right) y,\left( {1 - \lambda }\right) \left( {5 - {3z}}\right) }\right)\) ; (3) \(\lambda  = 1\) ,

势函数为 \(\frac{1}{3}{x}^{3} + {5xy} - {2y} + {3xyz} - 2{z}^{2} + C\) .

5. \(\frac{4}{3}\) .

\section{* 第二十三章 向量函数微分学}

\subsection{习题 23.2}

2. (1) \({\mathbf{f}}^{\prime }\left( {{x}_{1},{x}_{2}}\right)  = \left( \begin{matrix} \sin {x}_{2} & {x}_{1}\cos {x}_{2} \\  2\left( {{x}_{1} - {x}_{2}}\right) &  - 2\left( {{x}_{1} - {x}_{2}}\right) \\  0 & 4{x}_{2} \end{matrix}\right) ,{\mathbf{f}}^{\prime }\left( {0,\frac{\pi }{2}}\right)  = \left( \begin{matrix} 1 & 0 \\   - \pi & \pi \\  0 & {2\pi } \end{matrix}\right)\) ;

(2) \({f}^{\prime }\left( {{x}_{1},{x}_{2},{x}_{3}}\right)  = \left( \begin{matrix} 2{x}_{1} & 1 & 0 \\  {x}_{2}{\mathrm{e}}^{{x}_{1} + {x}_{3}} & {\mathrm{e}}^{{x}_{1} + {x}_{3}} & {x}_{2}{\mathrm{e}}^{{x}_{1} + {x}_{3}} \end{matrix}\right) ,{f}^{\prime }\left( {1,0,1}\right)  = \left( \begin{matrix} 2 & 1 & 0 \\  0 & {\mathrm{e}}^{2} & 0 \end{matrix}\right)\) .

4. (1) \(\cos x + \sin x\) ; (2) \(\left( \begin{matrix} \cos \left( {{x}_{1} - {x}_{2}}\right) &  - \cos \left( {{x}_{1} - {x}_{2}}\right) \\   - \sin \left( {{x}_{1} - {x}_{2}}\right) & \sin \left( {{x}_{1} - {x}_{2}}\right)  \end{matrix}\right)\) ;

(3) \(\left( \begin{matrix} {x}_{2}^{2} - 2{x}_{1}{x}_{2} & 2{x}_{1}{x}_{2} - {x}_{1}^{2} \\   - 1 - {x}_{2} & 1 - {x}_{1} \end{matrix}\right) ;\;\left( 4\right) \left( \begin{matrix} 2{x}_{1}{x}_{2}^{2} & 2{x}_{1}^{2}{x}_{2} \\   - 2 & 2 \\   - 1 & 1 \end{matrix}\right)\) ;

(5) \(\left( \begin{matrix} 4{x}_{1}{x}_{2}^{2} + {16}{x}_{1}{x}_{2} & 4{x}_{1}^{2}{x}_{2} + 8{x}_{1}^{2} \\  2{x}_{1} & 3 \end{matrix}\right)\) ; (6) \(\left( \begin{matrix} 2{x}_{1}{x}_{2}^{2}{x}_{3}^{2} & 2{x}_{1}^{2}{x}_{2}{x}_{3}^{2} & 2{x}_{1}^{2}{x}_{2}^{2}{x}_{3} \\  2 & 2 & 2 \\  1 & 1 & 1 \end{matrix}\right)\) .

5. \(\left( {{h}_{x} + {h}_{u}{f}_{x} + {h}_{v}\left( {{g}_{x} + {g}_{u}{f}_{x}}\right) ,{h}_{u}{f}_{y} + {h}_{v}\left( {{g}_{y} + {g}_{u}{f}_{y}}\right) }\right)\) .

7. (1) \(\mathbf{f}\left( \mathbf{x}\right)  = \mathbf{x} + \mathbf{b};\;\left( 2\right) \mathbf{f}\left( \mathbf{x}\right)  = {\left( \int {\varphi }_{1}\left( {x}_{1}\right) \mathrm{d}{x}_{1},\cdots ,\int {\varphi }_{n}\left( {x}_{n}\right) \mathrm{d}{x}_{n}\right) }^{\mathrm{T}}\) .

8. (1) 极小值点 \({x}_{0} = {\left( -\frac{17}{6}, - \frac{7}{3}, - \frac{2}{3}\right) }^{\mathrm{T}}\) ; (2) 稳定点 \({x}_{0} = {\left( 0,0,0\right) }^{\mathrm{T}}\) ,非极值点.

9. (1) \({\left( f \circ  \mathbf{h}\right) }^{\prime } = \left( {{x}_{2} + 1,{x}_{1} - 1}\right) ,{\left( f \circ  \mathbf{t}\right) }^{\prime } = \left( {{x}_{2}{x}_{3} - 1,{x}_{1}{x}_{3} - 1,{x}_{1}{x}_{2} - 1}\right)\) ,

\({\left( \mathbf{h} \circ  \mathbf{g}\right) }^{\prime } = \left( {\cos x + \sin x}\right) \left( \begin{matrix} \cos x - \sin x \\   - 1 \end{matrix}\right) ,\)

\({\left( \mathbf{h} \circ  \mathbf{t}\right) }^{\prime } = \left( \begin{matrix} 2{x}_{1}{x}_{2}{x}_{3} + {x}_{2}^{2}{x}_{3} + {x}_{2}{x}_{3}^{2} & {x}_{1}^{2}{x}_{3} + 2{x}_{1}{x}_{2}{x}_{3} + {x}_{1}{x}_{3}^{2} & {x}_{1}^{2}{x}_{2} + {x}_{1}{x}_{2}^{2} + 2{x}_{1}{x}_{2}{x}_{3} \\   - {x}_{2}{x}_{3} + 1 &  - {x}_{1}{x}_{3} + 1 &  - {x}_{1}{x}_{2} + 1 \end{matrix}\right) ,\)

\({\left( \mathbf{s} \circ  \mathbf{g}\right) }^{\prime } = \left( \begin{matrix} \sin {2x} \\   - 2\sin x \\   - \sin x \end{matrix}\right) ;\;\left( 2\right) {\left( \mathbf{g} \circ  f \circ  \mathbf{h}\right) }^{\prime } = \left( \begin{matrix} \left( {{x}_{2} + 1}\right) \cos w & \left( {{x}_{1} - 1}\right) \cos w \\   - \left( {{x}_{2} + 1}\right) \sin w &  - \left( {{x}_{1} - 1}\right) \sin w \end{matrix}\right) ,\)

其中 \(w = {x}_{1}{x}_{2} - {x}_{2} + {x}_{1}\) ,

\[
{\left( \mathbf{s} \circ  \mathbf{t} \circ  \mathbf{s}\right) }^{\prime } = \left( \begin{matrix} {16}{x}_{1}^{3}{x}_{2}^{2}{\left( {x}_{2} + 4\right) }^{2} & {16}{x}_{1}^{4}{x}_{2}\left( {{x}_{2} + 2}\right) \left( {{x}_{2} + 4}\right) \\  4{x}_{1} & 6 \\  2{x}_{1} & 3 \end{matrix}\right) .
\]

12. (1) 凡使 \({\mathbf{a}}^{\mathrm{T}}\mathbf{x} = {t}_{0}\) ( \({t}_{0}\) 为 \(\varphi\) 的稳定点) 的解 \({\mathbf{x}}_{0}\) 均为 \(f\) 的稳定点. 习题 23.3

2. \({z}_{xx} = \frac{2xy}{{\left( {x}^{2} + {y}^{2}\right) }^{2}},{z}_{xy} = \frac{{y}^{2} - {x}^{2}}{{\left( {x}^{2} + {y}^{2}\right) }^{2}},{z}_{yy} = \frac{-{2xy}}{{\left( {x}^{2} + {y}^{2}\right) }^{2}}\) .

3. (3) \(\frac{\partial u}{\partial x} = \frac{1}{\Delta }\left( {{f}_{1}^{\prime }{g}_{3}^{\prime } - {f}_{3}^{\prime }{g}_{1}^{\prime }}\right) ,\frac{\partial u}{\partial y} = \frac{1}{\Delta }\left( {{f}_{2}^{\prime }{g}_{3}^{\prime } - {f}_{3}^{\prime }{g}_{2}^{\prime }}\right) ,\frac{\partial u}{\partial v} =  - \frac{u}{\Delta }{g}_{3}^{\prime }\left( {{f}_{1}^{\prime } + {f}_{2}^{\prime } + {f}_{3}^{\prime }}\right)\) ,其

中 \(\Delta  = {g}_{3}^{\prime }\left\lbrack  {1 + v\left( {{f}_{1}^{\prime } + {f}_{2}^{\prime } + {f}_{3}^{\prime }}\right) }\right\rbrack\) .

5. (1) \(\left( \begin{array}{ll} \frac{\partial x}{\partial u} & \frac{\partial x}{\partial v} \\  \frac{\partial y}{\partial u} & \frac{\partial y}{\partial v} \end{array}\right)  = \frac{1}{w}\left( \begin{matrix} u & v \\  u\arctan \frac{v}{u} - v & v\arctan \frac{v}{u} + u \end{matrix}\right)\) ,其中 \(w = \sqrt{{u}^{2} + {v}^{2}}\) ;

(2) \(\left( \begin{array}{ll} \frac{\partial x}{\partial u} & \frac{\partial x}{\partial v} \\  \frac{\partial y}{\partial u} & \frac{\partial y}{\partial v} \end{array}\right)  = \frac{1}{x\left( {{\mathrm{e}}^{x}\sin y - {\mathrm{e}}^{x}\cos y + 1}\right) }\left( \begin{matrix} x\sin y &  - x\cos y \\  \cos y - {\mathrm{e}}^{x} & \sin y + {\mathrm{e}}^{x} \end{matrix}\right)\) .

\subsection{第二十三章总练习题}

7. (2) \(c = \left\{  \begin{matrix} \operatorname{arccot}\frac{{\beta }_{1}}{{\beta }_{2}}, & {\beta }_{2} \neq  0, \\  \pi , & {\beta }_{2} = 0. \end{matrix}\right.\)
